\chapter{Traumatologische Notfälle -- \gAasRsN{RS~VII}}

\emph{Maintainer: Josef Emhofer}

\minitoc

\paragraph{Trauma}

\gDefinition{Als Trauma  bezeichnet man eine
{\textperiodcentered}\,Sch\"adigung,
{\textperiodcentered}\,Verletzung, oder
{\textperiodcentered}\,Wunde, die durch eine
Gewalteinwirkung physikalischer (mechanischer, thermischer) oder chemischer Art verursacht wurde.}

Die Untersuchung des traumatologischen Patienten
(Traumacheck) wird unter \ref{topic:traumacheck} erklärt. 

%Die Traumata lassen sich grob ind folgende Klassen Einteilen:
% 
% \begin{itemize}
%     \item Sch\"adel-Hirn-Trauma (SHT)
%     \item Thorax trauma
%     \item Abdominal trauma / Bauchtrauma
%     \item Beckentrauma
%     \item Verletzungen der Wirbels\"aule
%     \item Polytrauma
%     \item Verletzungen der Extremit\"aten
%     \item Verbrennungen
% \end{itemize}



%%%%%%%%%%%%%%%%%%%%%%%%%%%%%%%%%%%%%%%%%%%%%%%%%%%%%%%%%%%
%%%%%%%%%%%%%%%%%%%%%%%%%%%%%%%%%%%%%%%%%%%%%%%%%%%%%%%%%%%
%%%%%%%%%%%%%%%%%%%%%%%%%%%%%%%%%%%%%%%%%%%%%%%%%%%%%%%%%%%

\section{Einsatztaktische Überlegungen bei verunfallten Patienten}

\subsection{Unfallort}
Bevor die Patientenversorgung beginnen kann, müssen am Unfallort einige Vorkehrungen für den weiteren Einsatzablauf getroffen werden. Die wichtigsten Aufgaben sind:
\begin{itemize}
	\item \gEs{Absichern} der Unfallstelle und Selbstschutz
	(Blaulicht, Warnblinkanlage, Uniformjacke anziehen, ev. Helm aufsetzen)
	\item Evt. Nachalarmierung weiterer Einsatzkräfte:
	weitere Rettungsmittel bei mehreren Verletzten, Feuerwehr, Polizei (Alarmierung ausschließlich über das Journal, Meldung der Anzahl an Verletzten)
	\item Bei mehreren Verletzten Triage (siehe ...)
	\item Eingeklemmte Personen werden im Fahrzeug stabilisiert bis die Feuerwehr die Verunfallten rettet; Helme tragen!
	\item Crash-Rettung nur bei Fremdgefährdung, HWS stabilisieren!
\end{itemize}

\subsection{Unfallmechanismus}
Der Unfallmechanismus gibt wertvolle Hinweise auf wahrscheinliche bzw. mögliche Verletzungen des Patienten. Der Rettungssanitäter kann durch eine systematische Vorgehensweise bei der Patientenversorgung, die eine genaue Betrachtung des Verletzungsmechanismus beinhaltet, mithelfen, die Schäden eines Traumas begrenzt zu halten. Die wichtigsten Fragen dabei sind immer:
\begin{itemize}
    \item \emph{Was ist passiert?}
    \item \emph{Wie ist der Patient verletzt worden?}
\end{itemize}

\paragraph{Verkehrsunfälle}

Bei Verkehrsunfällen müssen bei der Einschätzung des Verletzungsbildes drei Faktoren beachtet werden:
\begin{itemize}
    \item Deformierungsgrad des Fahrzeugs (Indiz für die involvierten Kräfte)
    \item Deformierung von Teilen der Fahrzeugkabine (Hinweis für Aufschlagpunkt
    des Körpers im Fahrzeug)
    \item Deformierung (Verletzungsmuster) des Patienten (Anzeichen dafür, welche Körperteile direkt aufgeprallt sind)
\end{itemize}

\subparagraph{Unfälle mit Fahrzeugen}

Für Unfälle mit Fahrzeugen gilt allgemein:

\begin{itemize}
    \item Die Geschwindigkeit des Fahrzeugs vor dem Unfall spielt eine wichtige Rolle. Je schneller das Fahrzeug unterwegs war, desto schwerer kann die Verletzung sein. \gZ{Dabei steigt die Bewegungsenergie des Fahrzeugs mit dem Quadrat der Geschwindigkeit an. Eine Verdoppelung der Geschwindigkeit hat daher eine Vervierfachung der Energie zur Folge, welche im Falle eines Unfalls frei wird und die Verformung verursacht.}
    \item Die Gewalteinwirkung auf den Patienten entspricht in ihrer Richtung i.d.R. der Einwirkung auf das Fahrzeug.
    \item Wird ein Insasse aus dem Fahrzeug geschleudert, steigert sich das Risiko schwerer Verletzungen (SHT, Wirbelsäule, Thorax, Abdomen etc.) um 300 \% (Böhmer et al. 2006). Auch bei den im Auto verbliebenen Insassen ist mit schweren Verletzungen zu rechnen.
    \item Wenn ein Insasse getötet oder eingeklemmt wird, ist bei den weiteren Insassen das Risiko schwerer Verletzungen erhöht.
    \item Knieverletzungen bei einem Autounfall können Hinweise auf Oberschenkel- und Beckenfrakturen sowie Hüftgelenksverrenkungen sein.
\end{itemize}

\subparagraph{Frontalzusammenstöße}

Bei Frontalzusammenstößen sollten folgende Merkmale
Alarmzeichen für schwere Verletzungen (Thoraxtrauma mit Herzkontusion und Pneumothorax, HWS-Trauma, Leber- und Milzruptur, Beckenfraktur) sein:
\begin{itemize}
    \item Knieanstoß am Armaturenbrett,
    \item Eingedrückte, leicht nach außen gewölbte Windschutzscheibe (ev. Spinnennetzmuster) durch Anstoß mit dem Kopf,
    \item sichtbare Veränderungen des Radstandes (z. B. Rad zur Seite geknickt, Achsenverschiebung um mehr als 30 cm am Unfallfahrzeug),
    \item Fahrzeugverformung um mehr als 50 cm.
\end{itemize}

\subparagraph{Seitenaufprall}

Beim Seitenaufprall sind typische schwere Verletzungen: HWS-Verrenkung, Thoraxtrauma, Verletzung der Aorta, seitenabhängig Leber- oder Milzruptur, Frakturen des Beckens. Bei einem Heckaufprall sollte man an HWS-Trauma (Schleudertrauma) denken und auch eine Frontalaufprallkomponente (Insasse wird nach vorne geworfen) mit einbeziehen (Lenkradkontusion).

\subparagraph{Fußgänger vs. Fahrzeug}

Wird ein \gEs{Fußgänger
von einem Fahrzeug (frontal) erfasst} ist generell mit schweren Verletzungen (vor allem Schädel-Hirn-Trauma, Wirbelsäule) zu rechnen. \gZ{Zur besseren Einschätzung kann man folgende Faustregel (nach Böhmer et al. 2006) verwenden:
\begin{itemize}
\item bis 50 km/h: Aufschlagen des Kopfes auf die Kühlerhaube
\item bis 70 km/h: Aufschlagen des Kopfes auf die Windschutzscheibe
\item mehr als 70 km/h: der Fußgänger wird über das Dach geworfen
\end{itemize}}
Lebensbedrohliche Verletzungen sind bereits möglich, wenn Radfahrer oder Fußgänger mit mehr als 30 km/h durch ein Fahrzeug erfasst werden. Erkennbar kann dies durch Bremsspuren, die bei trockener Fahrbahn länger als 10-20 m sind, sein.

\paragraph{Sturz aus großer Höhe}

Bei Stürzen aus großen Höhen (z. B. Leiter, Gerüst) ist stets
an HWS-, Wirbelsäulen- und SHT zu denken. Bei Sturz auf die
gestreckten Beine findet man typischerweise eine Fraktur der
Fersenbeine und/oder eine Stauchungsfraktur der Wirbelsäule
(Th\,12/L\,1), sowie Schienbeinplateau-Frakturen. Ab einer
Fallhöhe von 6 m ist von einer Polytraumatisierung
auszugehen.

\paragraph{Stich- und Schussverletzungen}
Traumen, die durch spitze Gegenstände (Stich- und
Pfählungsverletzungen) oder Schusswaffenmunition
(Schussverletzungen) hervorgerufen werden, werden
\gEs{penetrierende Verletzungen} genannt. Stichverletzungen
durch Messer, Scheren oder ähnliche Gegenstände werden
zumeist mit einer geringeren Geschwindigkeit zugefügt als
Schusswaffenmunition und haben daher eine geringere sekundäre
Verletzungsfolge als Schussverletzungen. Wie jeder andere
Fremdkörper auch, sollte ein noch im Körper des Patienten
befindliches Messer niemals präklinisch entfernt werden.
Stichwunden können oft sehr klein, schwach blutend und daher
harmlos aussehen. Um das Verletzungsausmaß richtig abschätzen
zu können, sollten folgende Aspekte beachtet werden:

\begin{itemize}
\item Welche Körperteile sind betroffen?
\item Wie viele Einstichwunden gibt es?
\item Wie lang war die Klinge?
\item Wie war der Einstichwinkel? (Welche Organe könnten noch betroffen sein?)
\item Wurde nach dem Einstich das Messer (der Gegenstand) im Körper des Patienten umgedreht?
\end{itemize}

\subsection{Zusammenfassung}
Aus allen einsatztaktischen Überlegungen ergibt sich ein Bündel an Erstmaßnahmen, welche Sie bei Unfällen durchführen müssen.

\gMassnahmen{
\begin{itemize}
\item Sichern Sie die Unfallstelle ab!
\item Verschaffen Sie sich einen Überblick (Wie viele Verletzte? ...) um ggf. weitere Einsatzkräfte nachzufordern!
\item Erheben Sie den Verletzungsmechanismus! Stellen Sie (sich) folgende Fragen:
	\begin{itemize}
	\item Aus welcher Richtung hat die Kraft auf den Patienten eingewirkt?
	\item Wie schnell war das Fahrzeug unterwegs? bzw. Aus welcher Höhe ist der Patient gefallen?
	\item War der Patient angegurtet?
	\item Welche besondere Verformungen sind am Fahrzeug sichtbar?
	\item Wurde der Patient aus dem Fahrzeug geschleudert?
	\end{itemize}
\end{itemize}
}


\section{Wunden}\label{topic:wunden}\index{Wunden}

\gDefinition{Eine Wunde ist eine Gewebeschädigung der Haut durch äußere physikalische Einwirkungen.}

\subsection{Einteilung von Wunden}

Wunden lassen sich nach verschiedenen Gesichtspunkten einteilen bzw. unterscheiden. 

\gParallel{
Eine Einteilung bezieht sich auf die \gE{physikalische
Verletzungsursache}:

\begin{enumerate}
	\item \gT{Mechanische Wunden}:
	\index{Wunden!mechanische}Entstehen durch spitze oder
	stumpfe Gewalteinwirkung auf den Körper (Stichwunde, Schürfwunde, Schusswunde, Bisswunde, ...)
	\item \gT{Thermische Wunden}: \index{Wunden!thermische}
	Entstehen durch Kälte- oder Hitzeeinwirkung auf den Körper (Verbrennungen, Stromeinwirkung, Erfrierung, ...)
	\item \gT{Chemische Wunden}:
	\index{Wunden!chemische}Entstehen durch Verätzungen mit
	Laugen oder Säuren
\end{enumerate}
}{
Eine weitere Unterteilung kann nach der \gE{Tiefe der
Schädigung} des Gewebes vorgenommen werden.

\begin{itemize}
	\item \gEs{Oberflächliche Wunden:} Leichte Verletzungen,
	welche max. bis zur Lederhaut reichen.
	\item \gEs{Tiefe Wunden:} Neben der Verletzung der Haut
	können hier auch Muskeln, Sehnen, größere Blutgefäße innere Organe und Knochen betroffen sein.
	\item \gEs{Penetrierende Wunden:} tiefe Wunden mit
	Eröffnung mindestens einer Körperhöhle (Schädel, Brust, Bauch)
% 	\item \gEs{Geschlossene Wunden:} Verletzungen der
% 	unteren Hautschichten ohne Eröffnung der Oberhaut. (Hämatom) % gestrichen
% lt. STEUER
\end{itemize}
}




Unterschiedliche \gE{Entstehungsmechanismen} erzeugen
unterschiedliche \gE{Wundarten} bzw. \gE{Wundformen}.

\begin{itemize}
	\item \gEs{Schnittwunde:} Glatte Wundränder; kann tieferliegende Nerven und Blutgefäße verletzen
	\item \gEs{Stichwunde:} Glatte Wundränder; meist kleine Wunde dafür aber tiefer
	\item \gEs{Risswunde:} Unregelmäßige Wundränder; meist nur Schädigung der Haut und
	keine tieferliegenden Nerven und Blutgefäßen
	\item \gEs{Quetschwunde:} Unregelmäßige Wundränder; es kann zum Aufplatzen der Haut kommen (Platzwunde)
	\item \gEs{Rissquetschwunde (RQW):} Kombination aus Riss- und Quetschwunde
	\item \gEs{Schürf- und Kratzwunde:} oberflächliche Wunde
	\item \gEs{Bisswunden:} Kombination aus Riss-, Stich- und Quetschwunde
	\item \gEs{Ablederung:} großflächiger Gewebedefekt, z.B. Skalpierung, große Lappenwunden
	\item \gEs{Amputation:} Abtrennung eines Körperteiles
	\item \gEs{Pfählung/Einspießung:} Eindringen und Verbleiben von spitzen aber z.T. unregelmäßigen Fremdkörpern
	\item \gEs{Hämatom:} Bluterguss; Ansammlung von Blut im Gewebe außerhalb der Blutgefäße
	\item \gEs{Schusswunde:} meist nach außen kleine Wunde
	mit großer Wirkung, ev. Ausschuss größer als Einschuss
\end{itemize}



\begin{tabularx}{\linewidth}[H]{XXX}
\\\toprule
\multicolumn{3}{c}{{\gTabHeading{Wunden}}}
\\\midrule
\includegraphics[width=\linewidth]{./images/wunde_rqw1.jpg}
\captionof{figure}{\gAuthNotesPicLegalOk{}{}Rissquetschwunde
vor der Versorgung im Spital.
\gSourcePicture{Hauer}}
&
\includegraphics[width=\linewidth]{./images/wunde_rqw2.jpg}
\captionof{figure}{\gAuthNotesPicLegalOk{}{}Rissquetschwunde
nach der Versorgung im Spital.
\gSourcePicture{Hauer}}
&
\includegraphics[width=\linewidth]{./images/pulsader_schnitt.jpg}
\captionof{figure}{\gAuthNotesPicLegalOk{}{}Schnittverletzung
oberhalb der "`Pulsader"'.
\gSourcePicture{Hauer}}
\\\midrule
\includegraphics[width=\linewidth]{./images/trauma/stichwunde-klein.jpg}
\captionof{figure}{\gAuthNotesPicLegalNotOk{}{}Glatte
Wundränder und eigentlich ganz unauffällig: Die Stichwunde.
\gSourcePicture{Hauer}} 
&
\includegraphics[width=\linewidth]{./images/pneumothorax_stichverletzung_diskret.jpg}
\captionof{figure}{\gAuthNotesPicLegalOk{}{}Stichverletzung:
unauffällig. \gCa{Dieser Patient ist
lebensgefährlich verletzt!}
(Pneumothorax)\gSourcePicture{Hauer}} &
\includegraphics[width=\linewidth]{./images/trauma/stichwunde-thorax-tief-01.jpg}
\captionof{figure}{\gAuthNotesPicLegalNotOk{}{}Stichwunde
mit eröffneter Brusthöhle.
\gSourcePicture{Hauer}} 
\\\bottomrule
\end{tabularx}


\begin{tabularx}{\linewidth}[H]{XXX}
\\\toprule
\multicolumn{3}{c}{{\gTabHeading{Fremdkörper}}}
\\\midrule
\includegraphics[width=\linewidth]{./images/trauma/fremdkoerper-stich-fixierung.jpg}
\captionof{figure}{\gAuthNotesPicLegalNotOk{}{}Fixierung
eines penetrierenden Fremdkörpers.
\gSourcePicture{Hauer}} 
&
\includegraphics[width=\linewidth]{./images/trauma/fremdkoerper-stich-messer-01.jpg}
\captionof{figure}{\gAuthNotesPicLegalNotOk{}{}Stich mit einem Küchenmesser. Das Messer verbleibt im Körper {\ldots}~
\gSourcePicture{Hauer}} 
&
\includegraphics[width=\linewidth]{./images/trauma/fremdkoerper-stich-messer-02.jpg}
\captionof{figure}{\gAuthNotesPicLegalNotOk{}{}{\ldots} bis in den Operationssaal.
\gSourcePicture{Hauer}} 
\\\bottomrule
\end{tabularx}

\begin{tabularx}{\linewidth}[H]{XXX}
\\\toprule
\multicolumn{3}{c}{{\gTabHeading{Frakturen}}}
\\\midrule
\includegraphics[width=\linewidth]{./images/trauma/fraktur-offen-finger.jpg}
\captionof{figure}{\gAuthNotesPicLegalNotOk{}{}Nicht immer
spektakulär: Offene Fraktur des Mittelfingers.
\gSourcePicture{Ben Stephenson, Cleveland, OH, USA.
Lizenz: CC-BY 2.0}} &
\includegraphics[width=\linewidth]{./images/trauma/fraktur-us-nativ.jpg}
\captionof{figure}{\gAuthNotesPicLegalNotOk{}{}Unterschenkelfraktur
nach Verkehrsunfall. 
\gSourcePicture{Hauer}} 
&
\includegraphics[width=\linewidth]{./images/trauma/fraktur-us-roe.jpg}
\captionof{figure}{\gAuthNotesPicLegalNotOk{}{}Röntgenbild
der Unterschenkelfraktur. 
\gSourcePicture{Hauer}} 
\\\bottomrule
\end{tabularx}

\begin{tabularx}{\linewidth}[H]{XXX}
\\\toprule
\multicolumn{3}{c}{{\gTabHeading{Verätzungen und
Verbrennungen}}} \\\midrule
%\includegraphics[width=\linewidth]{./images/trauma/fremdkoerper-stich-fixierung.jpg}
% \captionof{figure}{\gAuthNotesPicLegalNotOk{}{}Fixierung
% eines penetrierenden Fremdkörpers.
% \gSourcePicture{Hauer}} 
&
\includegraphics[width=\linewidth]{./images/trauma/burn-chem01.jpg}
\captionof{figure}{\gAuthNotesPicLegalNotOk{}{}Verätzung.
\gSourcePicture{Hauer}} 
&
\includegraphics[width=\linewidth]{./images/trauma/burn-chem02.jpg}
\captionof{figure}{\gAuthNotesPicLegalNotOk{}{}Verätzung.
\gSourcePicture{Hauer}} 
\\\midrule
\includegraphics[width=\linewidth]{./images/trauma/combustio-sa-01.jpg}
\captionof{figure}{\gAuthNotesPicLegalNotOk{}{}Großflächige
Verbrennung. \gSourcePicture{Hauer}} 
&
\includegraphics[width=\linewidth]{./images/trauma/combustio-sa-02-hand.jpg}
\captionof{figure}{\gAuthNotesPicLegalNotOk{}{}Ablederung.
\gSourcePicture{Hauer}} 
&
\includegraphics[width=\linewidth]{./images/trauma/combustio-sa-03-escharotomie.jpg}
\captionof{figure}{\gAuthNotesPicLegalNotOk{}{}Escharatomie. Durch die Ringförmige Verbrennung am Thorax zieht sich die
Haut zusammen und die Atmung erschwert. Durch die
"`Entlastungsschnitte"' kann der Pateint wieder atmen. 
\gSourcePicture{Hauer}} 
\\\bottomrule
\end{tabularx}



\subsection{Wundversorgung}

Die unterschiedlichen Wundarten benötigen auch unterschiedliche
Wundversorgungen. Bei den mechanischen Wunden wird die Blutstillung im Vordergrund stehen. Bei
thermischen Wunden wie z.B. Verbrennungen ist die Kühlung der betroffenen Stelle
als erstes durchzuführen. Bei chemischen Wunden sollte die Verdünnung der
ätzenden Substanz mit Wasser primär erfolgen. Schlussendlich werden aber alle
Wunden mit einer keimfreien Wundauflage zum Schutz vor Infektionen bedeckt.

Je nach Stärke der Blutung versorgen Sie Wunden mit einem
Wundverband, einem Druckverband oder in extremen Ausnahmesituationen durch eine Abbindung. Als
Erstmaßnahme bevor ein Verband zur Verfügung steht können Sie mit einer sterilen
Kompresse die Wunde abdrücken oder auch eine zuführende Arterie abdrücken. Die
Techniken haben Sie im Erste Hilfe-Teil Ihrer Ausbildung bereits gelernt. An
dieser Stelle sollen nur die wichtigsten \gEs {Grundregeln
der Wundversorgung}\label{topic:wundversorgung-grundregeln} wiederholt werden:

\begin{itemize}
    \item Berühren Sie die Wunden nicht direkt, außer es ist
    zur raschen Blutstillung unbedingt notwendig!
    \item Arbeiten Sie immer mit Handschuhen und mit sterilem Material!
    %\item Wunden sollen grundsätzliche \gE{nicht} gereinigt
    %werden! Sie dürfen nur lose Glassplitter oder ähnliches
    %vorsichtig aus der Wunde entfernen.
    \item Wunden, die im Spital weiter versorgt werden, sollen
    nicht übermäßig gereinigt werden. \underline{Grobe}
    Verschmutzungen werden nur wenn unbedingt notwendig mit
    steriler Flüssigkeit \gZ{(Kochsalzlösung, Ringerlösung,
    o.ä.)} oberflächlich \underline{ab}gespült. Lose, kleine
    Fremdkörper wie z.B. Glassplitter dürfen entfernt werden.
    \item Nicht-weiterbehandlungspflichtige Wunden können
    mit einem nicht-alkoholischen
    Desinfektionsmittel\footnote{\gZ{z.B.
    Octenisept\texttrademark{}}} gereinigt werden.
    \item Belassen Sie Fremdkörper (Messer, Schere oder sonstige Pfählungsgegenstände) in der Wunde! Stabilisieren Sie diese Gegenstände mit geeignetem Fixationsmaterial (z.B. Mullbinden).
    \item Versuchen Sie \gE{nicht}, ausgetretene Strukturen
    (z.B. Gehirn, Eingeweide) zurück zu stopfen! Solche Organteile sollten mit steriler NaCl-Lösung feucht gehalten und anschließend steril mit einem Verbandstuch abgedeckt werden!
    \item Steriler Wundverband mit der jeweiligen Situation
    angemessenem Material (z.B. Pflaster, Momentverband,
    Kompresse mit Mullbinde, Wundfolie, \ldots)
\end{itemize}

\subsection{Gefahren bei Wunden}
Jede Wunde hat (in unterschiedlichem Ausmaß) eine Auswirkung auf den Gesamtorganismus. Je nach Wundart, stärke der Gewalteinwirkung, betroffener Struktur und Sorgfalt bei der Wundversorgung gehen verschiedene Gefahren von Wunden aus:
\begin{itemize}
    \item \gEs{Verletzung wichtiger Gewebestrukturen (Organe):} Je nach betroffener Körperregion treten unterschiedliche Verletzungsmuster auf. In den folgenden Abschnitten lernen Sie über spezielle traumatologische Notfälle wie z.B. Schädel-Hirn-Trauma, Throax-Trauma, Bauchtrauma etc.
    \item \gEs{Starke Blutung:} Bei großflächigen Wunden oder Verletzung wichtiger Blutgefäße kann es zu starken Blutungen und infolge dessen zu einem Volumenmangelschock kommen. Blutstillung und Schockbekämpfung sind die wichtigsten Maßnahmen.
    \item \gEs{Starke Schmerzen:} Treten im Zusammenhang mit einer Verletzung starke Schmerzen auf, muss ein Notarzt für eine Schmerztherapie gerufen werden.
    \item \gEs{Eintritt von Krankheitserregern
    (Infektionsgefahr):} Jede Wunde ist eine
    Eintrittsstelle für Krankheitserreger in unseren Organismus. Der Patient sollte daher über einen aufrechten Tetanus-Impfschutz verfügen (siehe Kapitel Hygiene)\footnote{Zur Verifizierung eines aufrechten Tetanus-Impfschutzes muss jede Wunde einem Arzt vorgestellt werden. Eine Impfung muss spätestens alle 10 Jahre aufgefrischt werden, bei entsprechenden Verletzungen schon früher.\cite{BoehmerTaschRett6}}. Bei Bisswunden kann über den Speichel einer erkrankten Tieres Tollwut übertragen werden. Versuchen Sie bei Bisswunden daher immer zu erfragen, wer der Besitzer des Tieres ist und vermerken Sie dies auf Ihrem Einsatzprotokoll. Unter Umständen kann es notwendig sein die Polizei einzuschalten.
\end{itemize}

%\subsection{Mechanische
%Wunden}\label{topic:wunden-meschanische}\index{Wunden!mechanische}





% \begin{tabularx}{\linewidth}[H]{XXX}
% \\\toprule
% \multicolumn{3}{c}{{\gTabHeading{Diashow Muster}}}
% \\\midrule
% \includegraphics[width=\linewidth]{./images/stop.jpg}
% \captionof{figure}{\gAuthNotesPicLegalNotOk{}{}Muster
% \gSourcePicture{}} 
% &
% \includegraphics[width=\linewidth]{./images/stop.jpg}
% \captionof{figure}{\gAuthNotesPicLegalNotOk{}{}Muster
% \gSourcePicture{}} 
% &
% \includegraphics[width=\linewidth]{./images/stop.jpg}
% \captionof{figure}{\gAuthNotesPicLegalNotOk{}{}Muster
% \gSourcePicture{}} 
% \\\midrule
% \includegraphics[width=\linewidth]{./images/stop.jpg}
% \captionof{figure}{\gAuthNotesPicLegalNotOk{}{}Muster
% \gSourcePicture{}} 
% &
% \includegraphics[width=\linewidth]{./images/stop.jpg}
% \captionof{figure}{\gAuthNotesPicLegalNotOk{}{}Muster
% \gSourcePicture{}} 
% &
% \includegraphics[width=\linewidth]{./images/stop.jpg}
% \captionof{figure}{\gAuthNotesPicLegalNotOk{}{}Muster
% \gSourcePicture{}} 
% \\\bottomrule
% \end{tabularx}

%\subsubsection{Stichwunden}\label{topic:stichwunden}\index{Stichwunden}

% \begin{gWideParEnv}
% \begin{tabularx}{\linewidth}[H]{XXX}
% \\\toprule
% \multicolumn{3}{c}{{\gTabHeading{Diashow}}}
% \\\midrule
%  &
% \includegraphics[width=\linewidth]{./images/stichverletzung_oberarm.jpg}
% \captionof{figure}{\gAuthNotesPicLegalOk{}{}Stichverletzung
% am Oberarm, Schwellung. \gSourcePicture{Hauer}} &
% {leer}
% \\\bottomrule
% \end{tabularx}
% \end{gWideParEnv}


%\subsubsection{Schnittwunden}\label{topic:schnittwunden}\index{Schnittwunden}



%\subsubsection{Rißwunden}\label{topic:risswunden}\index{Rißwunden}

%\subsubsection{Quetschwunden}\label{topic:quetschwunden}\index{Quetschwunden}

%\subsubsection{Rißquetschwunden
%(RQW)}\label{topic:rqw}\index{Rißquetschwunden}\index{RQW}

% \begin{gDiasEnv3}
% \includegraphics[width=\linewidth]{./images/wunde_rqw1.jpg}
% \captionof{figure}{\gAuthNotesPicLegalOk{}{}Rißquetschwunde
% vor der Versorgung im Spital.
% \gSourcePicture{Hauer}}
% &
% \includegraphics[width=\linewidth]{./images/wunde_rqw2.jpg}
% \captionof{figure}{\gAuthNotesPicLegalOk{}{}Rißquetschwunde
% nach der Versorgung im Spital.
% \gSourcePicture{Hauer}}
% &
% \end{gDiasEnv3}





%\subsubsection{Schußwunden}\label{topic:schusswunden}\index{Schußwunden}




%\subsection{Thermische Wunden}\label{topic:wunden-thermische}\index{Wunden!thermische}

%\subsubsection{Brandwunden}\label{topic:brandwunden}\index{Brandwunden}

%\subsubsection{Erfrierungen}\label{topic:erfrierungen}\index{Erfrierungen}

%\subsection{Chemische Wunden}\label{topic:wunden-chemische}\index{Wunden!chemische}

%\subsubsection{Verätzungen}\label{topic:veraetzungen}\index{Verätzungen}


%%%%%%%%%%%%%%%%%%%%%%%%%%%%%%%%%%%%%%%%%%%%%%%%%%%%%%%%%%%
%%%%%%%%%%%%%%%%%%%%%%%%%%%%%%%%%%%%%%%%%%%%%%%%%%%%%%%%%%%

%\subsection{Stichverletzungen}

%NIEMALS  HERAUSZIEHEN!!!!!!!!

%steril abdecken + fixieren

%verletzte Strukturen??, dient ev. als {\quotedblbase}St\"opsel{\textquotedblleft}??

%Fremdk\"orper soll nicht tiefer rutschen



%%%%%%%%%%%%%%%%%%%%%%%%%%%%%%%%%%%%%%%%%%%%%%%%%%%%%%%%%%%
%%%%%%%%%%%%%%%%%%%%%%%%%%%%%%%%%%%%%%%%%%%%%%%%%%%%%%%%%%%
%%%%%%%%%%%%%%%%%%%%%%%%%%%%%%%%%%%%%%%%%%%%%%%%%%%%%%%%%%%

\section{Schädel-Hirn-Trauma (SHT)}
%Folge von Gewalteinwirkung auf den  Sch\"adel
\gDefinition{Gewalteinwirkung auf den Kopf, die zusätzlich zu den nicht immer
vorhandenen Hautwunden und zu den Frakturen des knöchernen Schädels
Funktionsstörungen und Verletzungen des Gehirns hervorrufen.\cite{Gorgass7}} Die häufigsten Ursachen von Schädel-Hirn-Traumata (SHT) sind Verkehrsunfälle, Stürze und Gewaltverbrechen. Je nach Schwere und Mechanismus können verschiedene Formen des SHT auftreten.

\gFazit{Das Leitsymptom eines jeden SHT ist die
\gEs{Bewusstseinsstörung}.}

\paragraph{Arten}

\begin{itemize}
\item Gedecktes SHT: ohne  Er\"offnung der harten Hirnhaut (Dura  mater)
\item Offenes SHT: mit  Er\"offnung der harten Hirnhaut (Dura  mater)
\end{itemize}

\paragraph{Begleitende Verletzungen}

\begin{itemize}
    \item Verletzung der Kopfschwarte  (Ri{\ss}quetschwunde)
    \item Frakturen von
    \begin{itemize}
        \item Sch\"adeldach
        \item Sch\"adelbasis (Blutung bzw. Liquoraustritt
        aus Ohren oder Nase, Monokel-
        oder Brillenhämatom)
        \item Gesichtssch\"adel
        \item Unterkiefer
    \end{itemize}
\end{itemize}

\gCave{Die Schwere eines SHT lässt sich nur bedingt von den äußerlich sichtbaren Verletzungen ableiten! Oft werden im Krankenhaus Frakturen festgestellt, ohne dass äußere Wunden vorliegen. Eine sorgfältige Untersuchung ist daher bei jedem SHT Pflicht!}

\gMassnahmen{
    \subsubsection{Allgemeine Maßnahmen beim
    SHT}\label{topic:sht-massnahmen-allgemein}
    
    \begin{itemize}
        \item Vitale Bedrohung beurteilen. \par \gTxBeiVit
        \gTxMassVitBes
        \begin{itemize}
            \item \ce{O2} 10-15\gLMin
            \item Lagerung: s.u.
            \item Monitoring: Zusätzlich Beurteilung nach GCS (Verlauf beobachten und
        dokumentieren)\gSee{topic:gcs}
        \end{itemize}
        \item Lagerung: leicht erhöhter Oberkörper, ca.
        30\,\textdegree{}, wenn keine Kontraindikation
        besteht.
        \item je nach Unfallmechanismus HWS-Immobilisation (z.B. mit Stifneck\texttrademark)
        \item Ggfs. Wundversorgung
        \begin{itemize}
                \item steriler Verband (steriles Abdecken)
        \item Knochen/Fremdkörper nicht entfernen/nicht zurück stopfen        
    \end{itemize}
%        \begin{itemize}
%            \item Gedecktes SHT:
%            \begin{itemize}
%                \item Steriler Verband
%            \end{itemize}
%            \begin{itemize}
%                \item steriler Verband locker auf offene
%                Verletzungen
%                \item Knochen / Hirnteile NICHT 
%                zur\"uckstopfen /  entfernen!
%                \item Ausgetretene Hirnmasse feucht steril abdecken
%            \end{itemize}
%        \end{itemize}
    \end{itemize}
}

\subsection{Gedecktes SHT}

Ein gedecktes Schädel-Hirn-Trauma ist ein SHT ohne Eröffnung der harten
Hirnhaut. Je nach Schwere wird das SHT in drei Grade eingeteilt:

\begin{itemize}
\item \gT{Gehirnersch\"utterung} -- SHT 1{\textdegree} --
Commotio cerebri
\item \gT{Gehirnprellung} -- SHT 2{\textdegree} -- Contusio
cerebri
\item \gT{Gehirnquetschung} -- SHT 3{\textdegree} 
\end{itemize}

\subsubsection{Gehirnersch\"utterung (Commotio cerebri) --
SHT 1\textdegree}
\index{Gehirnersch\"utterung}\index{Commotio (cerebri)}

\gMarried{Die Gehirnerschütterung ist die leichteste Form des
Schädelhirntraumas. Leitsymptom ist ein Trauma mit
anschließender kurzer Bewusstlosigkeit. Weiters kann es zu
flüchtigen neurologischen Ausfällen kommen, es kann auch eine
retrograde Amnesie (Gedächtnisverlust über den Vorgang
bzw. Unfall) bestehen. Eventuell klagt der Patient über
Übelkeit erbrechen und Schwindel. Die Gehirnerschütterung
zieht keine bleibenden Schäden nach sich.}
{
\begin{itemize}
    \item  Bewu{\ss}tlosigkeit/Bewu{\ss}tseinstr\"ubung
    {\textless}\,1\,h
    \item \"Ubelkeit, Erbrechen, Schwindel
    \item evt. fl\"uchtige neurologische Ausf\"alle
    \item retrograde  Amnesie  ({\quotedblbase}kann sich an den
    Unfallhergang nicht erinnern{\textquotedblleft} )
    \begin{itemize}
        \item im CT/MR 
        keine nachweisbaren Sch\"aden
        \item keine  bleibenden Sch\"aden
    \end{itemize}
\end{itemize}}

%\gMassnahmen{
%    \paragraph{Spezielle Maßnahmen beim SHT ---Gehirnerschütterung}
%    \begin{itemize}
%        \item 
%    \end{itemize}
%}


\subsubsection{Gehirnprellung (Contusio cerebri) --
SHT 2\textdegree}

\index{Gehirnprellung}\index{Contusio (cerebri)}

\gMarried{Die Gehirnprellung ist eine schwerere Form des Schädel-Hirntraumas. 
Leitsymptom ist die \gEs{Bewusstlosigkeit größer
 als 1\,h}, beziehungsweise die Bewusstseinseintrübung
 >\,24\,h. Mittels weiterführender Diagnostik im Spital
 lassen sich die \gEs{Hirnschädigungen} nachweisen. Es kann zu
 Nervenverletzungen kommen, dies kann zu vegetativen Störungen führen.
 Bleibende Schäden sind möglich.}
 {
 \begin{itemize}
    \item Bewu{\ss}tlosigkeit  {\textgreater}\,1\,h
    \item Bewu{\ss}tseinstr\"ubung {\textgreater}\,24\,h
    \item nachweisbare Gehirnsch\"adigung / bleibende  Ausf\"alle
    \item ev. Gehirnnervenabrisse 
    \item  vegetative St\"orungen
\end{itemize}}

%\gMassnahmen{
%    \paragraph{Spezielle Maßnahmen beim SHT --- Gehirnprellung}
%    \begin{itemize}
%        \item 
%    \end{itemize}
%}




\subsubsection{Gehirnquetschung -- SHT 3{\textdegree}}

\index{Gehirnquetschung}

\gMarried{Die Gehirnquetschung ist eine sehr schwere Form des
Schädelhirntraumas. Es kommt zu einer \gEs{Bewusstlosigkeit
länger als 24\,h} und zu teils sehr ausgeprägten
Begleitsymptomen wie zum Beispiel Hirndruckzeichen, Atem und
Kreislaufstörungen bis hin zum Atemstillstand durch Einklemmen
des Atemzentrums. 

Es kann weiteres zu intrakraniellen
Blutungen kommen (epidural, subdural, subarachnoidal).
Außerdem kommt es zu verschiedenen neurologischen Ausfällen
und Symptomen wie zum Beispiel zerebrale
Kampfanfälle. Mit \gEs{bleibenden Schäden} und Langzeitfolgen ist
zu rechnen.}
{
\begin{itemize}
    \item Bewu{\ss}tlosigkeit  {\textgreater}\,24\,h
    \item Atmungs-/ Kreislaufst\"orungen
    \item Hirndruck symptomatik
    \item evt. Einklemmung des Atemzentrums ${\rightarrow}$
    zentraler Atemstillstand
    \item Behinderung der Hirndurchblutung
    \item Hirnblutungen
%     \begin{itemize}
%         \item epidural
%         \item subdural
%         \item subarachnoidal
%     \end{itemize}
    \item L\"ahmungen
    \item Zerebrale Kr\"ampfe
\end{itemize}}

%\gMassnahmen{
%    \paragraph{Spezielle Maßnahmen beim SHT --- Gehirnquetschung}
%    \begin{itemize}
%        \item 
%    \end{itemize}
%}

\gMassnahmen{
    \subsubsection{Spezielle Maßnahmen beim SHT --
    gedecktes SHT}
    \begin{itemize}
        \item Allgemeine Maßnahmen beim SHT
        (\gSee{topic:sht-massnahmen-allgemein})

    \end{itemize}
}


\subsection{Offenes SHT}

Hinweise auf ein offenes SHT sind:

\begin{itemize}
    \item gro{\ss}e RQW mit tast-/sichtbaren  Knochensplittern
    \item Hirnaustritt (schlechte Prognose)
    \item Blutungen aus Nase/Ohren
    \item Liquoraustritt, oft mit Blut vermischt
    ({\quotedblbase}Tupfertest{\textquotedblleft})
\end{itemize}

\gMassnahmen{
    \subsubsection{Spezielle Maßnahmen beim SHT -- offenes
    SHT}
    \begin{itemize}
        \item Allgemeine Maßnahmen beim SHT
        (\gSee{topic:sht-massnahmen-allgemein})
        \item Ausgetretene Hirnmasse feucht steril abdecken
        und locker befestigen
    \end{itemize}
}


\subsection{Intrakranielle Blutungen}

\paragraph{\gT{Luzides Intervall} (lichtes Intervall)}

Die Symptome können eventuell erst Stunden nach der
Verletzung einsetzen. Dazwischen ist der Patient unauffällig.
Diese Zeitspanne wird Luzides Intervall genannt und ist
sehr gefährlich da man den Patienten leicht falsch
einschätzen kann.

\gMarried{Intrakranielle Blutungen können sowohl in Folge eines SHT als auch spontan (z.B. durch Platzen eines Aneurysmas) auftreten. Aufgrund der inneren Blutung steigt der Druck im Gehirn an, im Extremfall kann es zu einer Hirnstammeinklemmung kommen. Patienten mit intrakranieller Blutung sind vital bedroht! Beim SHT ist zu beachten, dass die Symptome oft auch erst Stunden nach dem Unfall auftreten können - siehe luzides Intervall.}
{\begin{itemize}
    \item stärkste Kopfschmerzen (langsam oder plötzlich einsetzend)
    \item Bewusstseinsstörungen bis Bewusstlosigkeit
    \item Pupillendifferenz, verlangsamte Pupillenreaktion
    \item ev. Lähmungen oder Krämpfe
    \item ev. Atemnot, unregelmäßige Atmung und Kreislaufstörung
    \item ev. Druckpuls (Bradykardie mit hohen Blutdruckwerten als Folge des Hirndrucks)
    \item Übelkeit, Erbrechen
\end{itemize}}

%\gAuthNotes{Zusammenfassen, keine Unterteilung in Sub-,
%Epi- und Aubarachnoidalblutung, nur mehr "`Hirnblutung"'}


%\subsubsection{Epidurale Blutung} \index{Epidurale Blutung}

%\gDefinition{Blutung zwischen
%der harten Hirnhaut und dem Schädelknochen.}

%\gMarried{Bei der epiduralen Blutung kommt es zu einer Blutung zwischen
%der harten Hirnhaut und dem Schädelknochen, bedingt durch
%den Riss einer Hirnhaut-Arterie. Durch die Blutung steigt
%der Hrindruck, es kommt dabei zu \gE{Hirndruckzeichen}.}
%{\begin{itemize}
%    \item Ri{\ss} einer Meningealarterie
%    \item Einblutung zwischen Kalotte und Dura mater
%    \item {\textquotedblleft}lichtes{\textquotedblright} (luzides )
%    Intervall: einige Stunden
%\end{itemize}}



%\subsubsection{Subdurale  Blutung} \index{Subdurale Blutung}

%\gDefinition{Blutung
%unterhalb der harten Hirnhaut und oberhalt der weichen
%Hirnhaut.}

%\gMarried{Bei der subduralen Blutung kommend es zu einer ein Blutung
%unterhalb der harten Hirnhaut. Der Verlauf ist oft
%schleichend und kann sich über mehrere Tage erstrecken. Die
%genaue Unterscheidung zwischen epiduraler und subduraler
%Blutung ist präklinisch nicht möglich.}
%{\begin{itemize}
%    \item schleichender (!) Verlauf: oft \"uber Tage
%    \item auch ohne  Trauma: Hypertoniker, Alkoholiker
%\end{itemize}}



%\subsubsection{Subarachnoidale Blutung -- SAB}
%\index{Subarachnoidale Blutung} (\index{SAB}SAB)

%\emph{Abkz.: \gT{SAB}}

% TODO test

%\gDefinition{Bei der subarachnoidalen Blutung kommt es zu einer
%Blutung unterhalb der Spinnengewebeshaut (Arachnoidea).} 

%\gMarried{Dafür gibt es
%unterschiedliche Ursachen, sie kann sowohl durch ein Trauma
%ausgelöst werden als auch spontan entstehen. Bei der
%\gE{spontanen} Subarachnoidalenblutung platzt typischerweise
%ein bereits bestehendes kleines Aneurysma eines Hirngefäßes und
%verursacht die Blutung. Dies kann aus heiterem Himmel
%passieren.

%\gFazit{Die subarachnoidalen Blutung kann durch ein Trauma
%oder spontan entstehen. Sie kann ein traumatologischer oder
%auch neurologischer Notfall sein.}

%Leitsymptom ist der plötzliche, schwere,
%\gEs{"donnerschlagartige" Kopfschmerz}, welcher mit
%Hirndruckzeichen und einem oft raschen Verfall des
%Patienten einhergeht. Der Patient ist akut \gE{vital
%bedroht!}}
%{
%\begin{itemize}
%    \item Trauma oder spontan
%    \item z.B. geplatztes Aneurysma, Br\"uckenvenenabri{\ss} bei
%    Decelerationstrauma
%    \item Klassisch: pl\"otzlicher, schwerer,
%    {\quotedblbase}donnerschlagartiger{\textquotedblleft} Kopfschmerz
%    \item evt. Herd symptomatik, schneller
%    Verfall,
%\end{itemize}}





\subsection{Spezielle Diagnostik beim SHT}

Grundsätzlich ist beim Schädel-Hirn-Trauma die übliche
Diagnostik durchzuführen, die besonders beim SHT wichtigen
Punkte sind hier jedoch nochmal beschrieben.

Beurteilung von:
\begin{itemize}
	\item Vitalfunktionen
	\item Neurologie (Neurocheck)
	\item Verletzungen (Traumacheck)
	\item Unfallmechanismus (Anamnese)
\end{itemize}

\paragraph{Neurologische Beurteilung}

Die neurologische Beurteilung ist beim SHT natürlich
besonders wichtig und besonders genau durchzuführen.
Auffälligkeiten können Hinweise auf eine direkte
Hirnschädigung oder eine Hirnblutung sein. 

\gParallel{
\begin{description}
\item[Bewu{\ss}tseinslage] klar/getr\"ubt/bewu{\ss}tlos
\item[Motorik] 
\item[Pupillenweite] Hirndruckzeichen
\item[Pupillenreaktion] Hirndruckzeichen
\end{description}

}{
\includegraphics[width=\linewidth]{./images/trauma/pupillendifferenz.jpg}
\captionof{figure}{\gCa{Alarmzeichen:} Eine Pupillendifferenz kann ein wichtiger
Hinweis auf einen erhöhten Hirndruck sein!} }


\paragraph{Traumacheck} Ein Traumackeck ist obligat
und sollte keiner weiteren Erwähnung bedürfen (siehe
\ref{topic:traumacheck}).



\gAuthNotes{GABS: Die einzelnen Massnahmen in Bezug auf die
unterschiedlichen SHTs gehören definiert.}

%\gMassnahmen{
%    \paragraph{Spezielle Maßnahmen beim SHT --- Gehirnerschütterung}
%    \begin{itemize}
%        \item 
%    \end{itemize}
%}

\begin{itemize}
\item (Gefahr: hohe Querschnittl\"ahmung)
\item Sch\"adel/HWS/BWS bilden eine Linie
\item Jedes SHT ist im Spital abzukl\"aren! (klinisch und  bildgebend)
\end{itemize}

\gCave{\gEs{Jedes SHT ist im Spital abzukl\"aren!} Es
besteht die Gefahr dass sich der Patient gegenw\"artig in der
\gEs{"`luciden Intervall"'} 
befindet und sich sein Zustand erst in einigen Stunden, dann aber schlagartig verschlechtert. Im Zweifelsfall wird der
Patient dort zur  Beobachtung (24h lang) aufgenommen.}



%%%%%%%%%%%%%%%%%%%%%%%%%%%%%%%%%%%%%%%%%%%%%%%%%%%%%%%%%%%
%%%%%%%%%%%%%%%%%%%%%%%%%%%%%%%%%%%%%%%%%%%%%%%%%%%%%%%%%%%
%%%%%%%%%%%%%%%%%%%%%%%%%%%%%%%%%%%%%%%%%%%%%%%%%%%%%%%%%%%


\section{Wirbelsäulentrauma}

\gAuthNotes{KOCH: Überarbeiten, evt etwas kürzen.}

\gDefinition{Gewalteinwirkung auf die Wirbelsäule, die zur Verschiebung oder zur Fraktur von Wirbeln mit oder ohne Rückenmarkschädigung führt.\cite{Gorgass7}}

\gParallel{
\paragraph{Typische Unfallmechanismen}

\begin{itemize}
    \item Sturz aus gro{\ss}er H\"ohe (> 3 m)
    \item Motorradunfall
    \item eingeklemmte Personen bei VU
    \item Schlag auf Kopf (Stauchung der HWS)
    \item Schleuderbewegung bei VU, wenn Patient nicht angeschnallt war oder Kopfstütze fehlt
    \item Suizidversuch durch Erhängen
\end{itemize}
}{
\paragraph{Gefahren}

\begin{itemize}
    \item Kompression des R\"uckenmarks (aufgrund von Einblutungen ins  R\"uckenmark und dadurch bedingtes Ödem) 
    \item Durchtrennung des R\"uckenmarks
    \item L\"ahmungen (Parese)
    \item Sensibilit\"atsst\"orungen (Par\"asthesie)
\end{itemize}
}





\gMassnahmen{
    \paragraph{Spezielle Maßnahmen beim Wirbelsäulentrauma --- Allgemein}
    
    \begin{itemize}
        \item bei ansprechbaren Patienten: immer HWS-Immobilisation (Stifneck)
        \item Rettung mit Schaufeltrage oder KED-System
        \item Lagerung auf/Schienung mit Vakuummatratze (wenn Schocklagerung notwendig, Tragentisch schräg stellen)
        \item ev. Sauerstoffgabe
        \item bei Querschnittslähmung: Kennzeichnung der Sensibilitätsgrenze mit Angabe der Uhrzeit
        \item ev. Wundversorgung
        \item ev. Schockbekämpfung
        \item regelmäßige Kontrolle der Atmung
        \item Erheben und Dokumentieren des Unfallmechanismus
    \end{itemize}
}

\subsection{Schleudertrauma, Peitschenschlag}

Schleudertraumen kommen bei Auffahrunf\"allen vor, insbesondere dann, wenn der Patient nicht angeschnallt war, oder die Kopfstütze fehlt. Es kommt zu einer Zerrung des Halsmarks, welche sich durch
\begin{itemize}
	\item Schmerzen im HWS-Bereich
	\item Kopfschmerzen
	\item Schwindel, Tinnitus
	\item ev. neurologische Ausfälle (in schweren Fällen)
\end{itemize}
bemerkbar macht.

\subsection{Querschnittssyndrom}

\gMarried{Kommt es zu einer Druckerhöhung (z.B. aufgrund
einer Blutung) im Rückenmark oder zu einer Durchtrennung des
Rückenmarks können Symptome eines Querschnitts auftreten. Ein
sorgsamer und vorsichtiger Umgang bei der Rettung, der
Untersuchung und dem Transport des Patienten ist unbedingt
notwendig.}{
\begin{itemize}
	\item Ausfall von Motorik und/oder Sensibilität
	\item Blutdruckabfall aufgrund fehlender Gefäßregulation
	\item Schlaffe oder spastische Lähmung der Extremitäten
	\item Inkontinenz aufgrund von Blasen-/Darmlähmung
	\item Dauererektion (Priapismus)
\end{itemize}
}

\subsection{Spinaler Schock}

\gMarried{Aufgrund des Unfalls kommt es zu einer plötzlichen Durchtrennung des Rückenmarks, wodurch die Gefäßregulation (und somit der Gefäßtonus) versagt. Somit sinkt der Blutdruck rasch ab, es kommt zu einer lebensbedrohlichen Kreislaufstörung.}{
\begin{itemize}
	\item deutlicher Blutdruckabfall aufgrund fehlender Gefäßregulation
	\item Querschnittslähmung	
	\item Atemlähmungen (bei hohem Querschnitt (im Bereich der HWS))
\end{itemize}
}

%%%%%%%%%%%%%%%%%%%%%%%%%%%%%%%%%%%%%%%%%%%%%%%%%%%%%%%%%%%
%%%%%%%%%%%%%%%%%%%%%%%%%%%%%%%%%%%%%%%%%%%%%%%%%%%%%%%%%%%
%%%%%%%%%%%%%%%%%%%%%%%%%%%%%%%%%%%%%%%%%%%%%%%%%%%%%%%%%%%

\section{Thoraxtrauma}

\gDefinition{Verletzung des Brustkorbes und/oder der darin
enthaltenen Organe.}

Thoraxverletzunge müssen als besonders gefährlich eingestuft
werden, da im Brustkorb (Thorax) lebenswichtige Organe (Herz,
Lunge) und große Gefäße liegen.  Nahezu jeder 2. Verkehrstote
ist Opfer eines Thoraxtraumas. \cite{Gorgass7}

\gCave{Sie sollten niemals aufgrund der äußeren Verletzungen auf die Schwere des Thoraxtraumas schließen, da schwere innere Verletzungen auftreten können!}

\paragraph{Symptome des Thoraxtraumas}~

\begin{tabularx}{\linewidth}{XX}
\gTabHeading{Allgemeine Symptome}
&
\gTabHeading{Spezifische Symptome}
\\\midrule
\begin{itemize}
    \item Dyspnoe
    \begin{itemize}
        \item Tachypnoe (beschleunigte Atmung)
        \item Zyanose
    \end{itemize}
    \item atemabh\"angige  Schmerzen,  Inspirationsstop
    \item ev. Hautemphysem (Luftansammlung in der Haut, typisches
    {\quotedblbase}Knistern{\textquotedblleft})
    \item Prellmarken (z.B. Lenks\"aule von PKW)
    \begin{itemize}
        \item gr\"undlicher Traumacheck!
    \end{itemize}
\end{itemize}
&
\begin{itemize}
    \item instabiler Thorax
    \begin{itemize}
        \item Sternumfraktur
        \item Serienrippenfraktur (zus. paradoxe Atmung --  Brustkorb senkt sich beim
        Einatmen)
    \end{itemize}
    \item Dyspnoe + Prellmarken +  Bluthusten
    \begin{itemize}
        \item Lungenprellung: Gasaustausch, da  Lungenabschnitte
        gesch\"adigt
    \end{itemize}
    \item Brustschmerz, Schock
    \begin{itemize}
        \item Rhythmusst\"orungen ${\rightarrow}$ Verletzung des  Herzens
        \item schwerer Schock, schneller Verfall ${\rightarrow}$ Riss der  Aorta
        \item Blut im Schwall aus dem Mund ${\rightarrow}$
        Speiser\"ohrenverletzung
    \end{itemize}
    \end{itemize}
\\
\end{tabularx}



\subsection{Pneumothorax}\label{topic:pneumothorax}\index{Pneumothorax}

\gDefinition{Luftansammlung im Pleuraraum}

(Vgl. Kap. Anatomie, \ref{topic:atemmechanik} und
\ref{topic:pleura})
\index{Pneumothorax!geschlossener}\index{Pneumothorax!offener}\index{Pneumothorax!Spontan-}\index{Pneumothorax!Spannungs-}\index{Spontanpneumothorax}\index{Spannungspneumothorax}

Zwischen Lungenoberfläche und innerer Thoraxwand herrscht
beim Gesunden ein Unterdruck gegenüber dem äußeren Luftdruck.
Die Lunge ist dadurch flexibel im Brustkorb aufgespannt. Beim
Pneumothorax gelangt durch eine Verletzung Luft in diesen
Pleuraspalt, welche den Unterdruck aufhebt und die Lunge
zieht sich zusammen. Man unterscheidet zwischen \gT{offenem
Pneumothorax} (Luft strömt durch eine Verletzung der
Brustwand von außen in den Pleuraspalt) und
\gT{geschlossenem Pneumothorax} (Luft strömt durch eine
Verletzung der Lungenoberfläche von innen in den
Pleuraspalt). Entsteht durch die Form der Wunde ein
Ventilmechanismus wird beim Einatmen Luft in den Pleuraspalt
gesogen, während sich beim Ausatmen die Wunde verschließt und
die Luft nicht mehr entweichen kann. Dadurch füllt sich der
Pleuraspalt mehr und mehr mit Luft, bis schließlich im
betroffenen Pleuraspalt ein Überdruck entsteht. Im Extremfall
drückt der Überdruck den Mittelfellraum auf die Gegenseite,
wodurch auch noch die gesunde Lungenhälfte und ev. die
Hohlvene komprimiert wird. In diesem Fall spricht man vom
\gT{Spannungspneumothorax}.
Ein Pneumothorax kann auch ohne vorangegangene Verletzung auftreten. Diese Form des
Pneumothorax nennt man \gT{Spontanpneumothorax} und tritt
am häufigsten bei jungen Menschen auf. Symptome des
Spontanpneumothorax sind Atemnot und Todesangst.



%\begin{itemize}
%   \item Lungen- (innerer) oder  Brustwandverletzung (\"au{\ss}erer)
%    \begin{itemize}
%        \item Luft gelangt in den Pleuraspalt
%        \item Unterdruck aufgehoben
%        \item Lunge kollabiert auf betroffener Seite
%    \end{itemize}
%    \item "`unkomplizierter"' Pneumothorax
%    \begin{itemize}
%        \item Luft kann ungehindert in Pleuraspalt und  wieder heraus
%        \item {\quotedblbase}nur{\textquotedblleft} Ausfall des betroffenen
%        Lungenfl\"ugels
%    \end{itemize}
%    \item Ventil-, \index{Spannungspneumothorax}Spannungspneumothorax
%    \begin{itemize}
%        \item Luft kann zwar hinein, aber nicht wieder  heraus
%        \item Druck auf betroffener Seite steigt ${\rightarrow}$  Organe (Lunge,
%        Herz) werden auf die andere Seite gedr\"uckt
%        \begin{itemize}
%            \item gesunde Seite wird verdr\"angt + Abknickung  gro{\ss}er
%            Gef\"a{\ss}e!
%        \end{itemize}
%    \end{itemize}
%\end{itemize}


%Therapie

%Pneumothorax, unkompliziert:

\begin{table}

\begin{gWideParEnv}

\begin{tabularx}{\linewidth}[H]{XXX}
\\\toprule
\multicolumn{3}{c}{{\gTabHeading{Pneumothorax}}}
\\\midrule
\includegraphics[width=\linewidth]{./images/pneumothorax_xray_halb.jpg}
\captionof{figure}{\gAuthNotesPicLegalNotOk{}{}Pneumothorax. 
beachte die Randlinie zwischen Lunge
(unten, innen) und dem luftgeüllten
Pleuraspalt (oben, außen) \gSourcePicture{Hauer}} &
\includegraphics[width=\linewidth]{./images/pneumothorax_xray_endstadium.jpg}
\captionof{figure}{\gAuthNotesPicLegalNotOk{}{}Kompletter
Pneumothorax der rechten Seite. Die Lunge ist kollapiert.  
\gSourcePicture{Hauer}} 
&
\includegraphics[width=\linewidth]{./images/pneumothorax_stichverletzung_diskret.jpg}
\captionof{figure}{\gAuthNotesPicLegalOk{}{}Diese
unauffällige Stichverletzung hat einen Pneumothorax
verursacht.
\gCa{Dieser Patient ist lebensgefährlich verletzt!} 
\gSourcePicture{Hauer}} \\\bottomrule
\end{tabularx}
\end{gWideParEnv}
\end{table}

\gMassnahmen{
    \paragraph{Spezielle Maßnahmen --- Pneumothorax}
    \begin{itemize}
        \item \gTxMassVitBes
        \begin{itemize}
            \item \ce{O2}: 10\,-\,15\,l
            \item Lagerung: Oberkörper hoch. Wenn stabile
            Seitenlage notwendig ist dann Lagerung auf die
            \gE{verletzte} Seite, damit die unverletzte
            Lunge  nicht unnötig belastet wird.
        \end{itemize}
        \item Bei offener Thoraxverletzung: \gEs{Nicht
        luftdicht verbinden!} Es kann sonst ein
        Spannungspneumothorax entstehen. Steril abdecken. 
    \end{itemize}
}

\gZ{Der Notarzt kann mittels Drainage einen Pneumothorax
entlasten.}

%\gAuthNotes{KOCH: 

%Es fehlt:
%Serienrippenfraktur erklären
%Paradoxe Atmung
%Gefahr von Blutverlust
%Therapie
%Ev. auch kurz etwas über STernumfraktur schreiben?}

\subsection{Rippenfraktur, Serienrippenfraktur und instabiler Thorax}\label{topic:serienrippenfraktur}\index{Serienrippenfraktur}
\gMarried{Von einer \gEs{Serienrippenfraktur} spricht man,
wenn mindestens drei benachbarte Rippen gebrochen sind. Sind
die Rippen an verschiedenen Stellen gebrochen oder gar
zertrümmert, werden diese Teile bei der Einatmung nach innen
gezogen und bei der Ausatmung nach außen gedrückt. Man
spricht in diesem Fall von \gEs{paradoxer Atmung} bzw.
\gEs{instabilem Thorax}, welche eine starke Beeinträchtigung
der Lungenfunktion nach sich zieht. Die Folge ist hochgradige
Atemnot. Außerdem treten (z.T. heftige) Schmerzen auf, die
sich bei der Einatmung verstärken. Rippenbruchstücke können
darunter liegende Organe verletzen. Je nach Rippe(n) kann es
zu Lungeneinrissen, Leber- und Milzverletzungen mit starken
inneren Blutungen kommen. Der Blutverlust kann 2-3 Liter
betragen. \cite{LPN3} In diesem Fall herrscht höchste
Schockgefahr.}{
\begin{itemize}
	\item (z.T. schwere) Atemnot, Zyanose, Tachypnoe
	\item (z.T. heftige) Schmerzen beim Einatmen
	\item ev. Schocksymptome
	\item ev. paradoxe Atmung
\end{itemize}
} 

\subsection{Sternumfraktur}\label{topic:sternumfraktur}\index{Sternumfraktur}
Eine Fraktur des Brustbeins (Sternum) tritt vor allem bei nicht angeschnallten Fahrzeuglenkern auf, wenn der Brustkorb gegen das Lenkrad prallt. Meistens treten starke Schmerzen auf, gelegentlich lassen sich eine Delle oder eine Stufe tasten.

\subsection{Verletzung des Herzens und der großen Gefäße}
Verletzungen des Herzens und der großen Gefäße treten am wahrscheinlichsten bei Stürzen aus großer Höhe und Verkehrsunfällen mit Frontalaufprall und großer Geschwindigkeit auf, seltener bei Schuss- oder Stichverletzungen. Folgende Verletzungen sind typisch:
\begin{itemize}
	\item \gEs{Aortenruptur}: Durch einen Einriss der Aorta kommt es zu schweren inneren Blutungen.
	\item \gEs{Herzkontusion}: Hier kommt es durch Quetschungen der Herzmuskulatur zu Sickerblutungen. Bei Sternumfrakturen mit Herzrhythmusstörungen sollte man an eine Herzkontusion denken.
	\item \gEs{Herzbeuteltamponade}: Es kommt zu einer Einblutung in den Herzbeutel, wodurch das Herz komprimiert wird und sich nicht mehr mit Blut füllen kann. Leitsymptome sind gestaute Halsvenen und schwerer Schockzustand (kardiogener Schock).
\end{itemize}

\section{Bauchtrauma}
\index{Bauchtrauma}

Ein Bauchtrauma oder Abdominaltrauma ist eine Verletzung des Abdomens. In einigen Büchern ist das Bauchtrauma als Untergruppe des akuten Abdomens klassifiziert. Häufigste Ursachen sind Verkehrsunfälle, Stürze aus großer Höhe sowie Stoß- und Schlagverletzungen. \cite{RedelsteinerHRNS1} Grundsätzlich ist zwischen offenem (perforierendem, penetrierendem) und geschlossenem (stumpfem) Bauchtrauma zu unterscheiden.

\begin{tabularx}{\linewidth}{XX}
\gTabHeading{offenes (perforierendes) Bauchtrauma}
&
\gTabHeading{geschlossenes (stumpfes) Bauchtrauma}
\\\midrule
    \begin{itemize}
        \item  Er\"offnung der Bauchh\"ohle, ev. Austritt von Darmschlingen
        \item Stich-, Pf\"ahlungsverletzung
        \item Schnittverletzung, {\quotedblbase}Platzbauch{\textquotedblleft}
    \end{itemize}
&
    \begin{itemize}
        \item Schlag, Tritt
        \item Prellung durch Lenkstange bei PKW-Unfall
    \end{itemize}
\\
\end{tabularx}



\subsection{Offenes Bauchtrauma}

\gMarried{
Offene Bauchtraumata sind seltener als geschlossene
Bauchtraumata. Beim offenen Bauchtrauma kann es zum Austritt
von Darmschlingen kommen. Die Größe der Eintrittswunde sagt
nichts über den Schweregrad der Verletzung aus, da unter der
Eintrittswunde unterschiedlich schwere Organschäden mit
inneren Blutungen versteckt sein können.
}{
    \includegraphics[width=\linewidth]{./images/trauma/bauchtrauma-offenes-01.jpg} \captionof{figure}{\gAuthNotesPicLegalOk{}{}Offenes Bauchtrauma mit Austritt von Darmschlingen.
    \gSourcePicture{Hauer}}
}


\gMassnahmen{
    
    \paragraph{Spezielle Maßnahmen beim offenen Bauchtrauma}
    \begin{itemize}
        \item \gTxMassVitBes
        \begin{itemize}
            \item Lagerung: Bauchdecke entspannen , Knierolle
        \end{itemize}
        \item Darmteile nicht  zur\"uckstopfen ${\rightarrow}$ sonst
        Darmri{\ss}/Darmverschlu{\ss}
        \item steril abdecken, mit Ringerl\"osung feuchthalten
        \begin{itemize}
            \item Gefahr des Eindringens von Keimen  (Peritonitis/Sepsis!)
        \end{itemize}
    \end{itemize}
}



\subsection{Geschlossenes Bauchtrauma}

Geschlossene Bauchtraumata sind aufgrund der inneren Blutungen besonders gefährlich. So kann es z.B. bei einer Milzruptur zu einem Blutverlust von 4 Liter kommen!

\gFazit{Immer nach \gEs{Prellmarken} suchen. \gEs{Patienten
entkleiden!}}

\gParallel{
\paragraph{Symptome}

\begin{itemize}
    \item starke  Bauchschmerzen
    \item brettharte  Bauchdecke  (Abwehrspannung, D\'efense )
    \item Leitsymptom: Schockzeichen
    \item beim Bewu{\ss}tlosen schwierig zu  diagnostizieren
    \begin{itemize}
        \item gr\"undlich nach \index{Prellmarken}\gEs{Prellmarken}
        suchen!
    \end{itemize}
\end{itemize}
}{
\paragraph{Gefahren}

\begin{itemize}
    \item Gef\"a{\ss}verletzungen
    \begin{itemize}
        \item Aorta, Vena cava,...
        \item massive innere Blutungen ${\rightarrow}$ Volumenmangelschock
    \end{itemize}
    \item Zerrei{\ss}ung  von Organen (\gEs{Milz},
    \gEs{Leber})
    \begin{itemize}
        \item massive \gEs{innere Blutungen} (bis über 4 Liter) ${\rightarrow}$
        Volumenmangelschock
    \end{itemize}
    \item Zerrei{\ss}ung  von Hohlorganen (Magen, Darm,Galleng\"ange )
    \begin{itemize}
        \item Austritt von Speisebrei, Magensaft, Stuhl, Gallensekret
        \begin{itemize}
            \item Bauchfellentz\"undung (Peritonitis) nach Std. bis  Tagen
            \item Sepsis (hohe Lethalit\"at), septischer Schock
        \end{itemize}
    \end{itemize}
\end{itemize}
}







\subsubsection{Milzruptur}\index{Milzruptur}

\gMarried{Die Milz ist ein teuflisches Organ. Sie besitzt rund um
das gut durchblutete Gewebe eine Bindegewebs-Kapsel. Das
Schlimme daran: Das Gewebe kann reißen, die Kapsel bleibt
jedoch intakt. Es kommt zuerst nur zu einer geringfügigen
Einblutung, dem Patienten geht es zunächst gut. Nach
einiger Zeit reißt dann die Kapsel und es kommt zu einer
ungehemmten Blutung in den Bauchraum. \gEs{Der Patient
verfällt von einem Moment auf den anderen in einen Volumenmangelschock.}}{
\begin{itemize}
	\item Schmerzen im linken Oberbauch
	\item manchmal Ausstrahlung der Schmerzen in die linke Schulter (Kehr-Zeichen)
	\item Volumenmangelschock
\end{itemize}
}

\paragraph{Primäre (einzeitige) Milzruptur} Die Milz und die Kapsel
reißen gleichzeitig, es kommt zu einem unmittelbar merkbaren
Schockzustand.

\paragraph{Sekundäre (zweizeitige) Milzruptur} Die Kapsel reißt erst
einige Zeit (Stunden bis Tage) nach dem Riss des Milzgewebes
\gE{(Latenzzeit)}. Es kommt zu einem \gEs{zeitlich versetztem} und oft
\gEs{unerwartetem Schockgeschehen}. Mitunter reißt die
Kapsel erst Stunden nach dem eigentlichen Unfall. \gEs{Daher
ist es wichtig jeden Patienten, bei dem auch nur der geringste
begründete Verdacht auf eine Milzverletzung besteht, einer
weiterführenden Untersuchung im Spital zuzuführen!}
\gZ{Mittels Ultraschall kann auf eine Milzblutung untersucht
werden.}

\gCave{Cave! {\quotedblbase}\index{Zweizeitige
Milzruptur}Zweizeitige Milzruptur{\textquotedblleft}}

\paragraph{Zusammenfassung:}
\begin{itemize}
    \item prim\"are Milzruptur: sofort, Kapsel durchgerissen
    \item sekund\"are Milzruptur:
    	\begin{itemize}
		\item Bei Unfall zunächst nur Riss des Milzgewebes, Kapsel noch intakt.
		\item symptomfreies Intervall (Latenzzeit), Sickerblutung
		\item Später reißt auch die Kapsel  ${\rightarrow}$ Schock! 
		\end{itemize} 
		%Sickerblutung, Kapsel anfangs noch intakt
    ${\rightarrow}$ Latenz bis Ri{\ss} der Kapsel
\end{itemize}




\subsubsection{Weitere innere Verletzungen}

Der Bauchraum beherbergt viele, teilweise sehr gut
durchblutete Organe die bei einer Verletzung zu erheblichen
Blutverlusten führen können. Präklinisch äußern sich diese
Verletzungen durch den Blutverlust hervorgerufenen
Schockzustand. Weitere hinweiszeichen sind Prellmarken oder
Schmerzen in der entsprechenden Gegend.

Grundsätzlich kann man ursächlich gegen diese Verletzung
wenig tun, im Vordergrund steht die Erkennung,
Schockbehandlung und der zügige Transport in ein
Traumazentrum.

\paragraph{Leberruptur} ~

\gMarried{Die Leber ist auch gut durchblutet und kann im Falle einer Verletzung große Probleme machen. Ursachen sind meist stumpfe Verletzungen, z.B. ein Lenkradaufprall oder eine Sicherheitsgurtkompression}{
\begin{itemize}
	\item Druckschmerz im rechten Oberbauch
	\item Volumenmangelschock
	\item späterer Kapselriss möglich (zweizeitige Ruptur wie bei Milz, siehe oben)
\end{itemize}
}

\paragraph{Aortenruptur}\index{Aortenruptur}\label{topic:aortenruptur}Durch
einwirkende starke Kräfte kann es zu einem Riss in der Aorta
kommen. Besonders bei Stürzen aus großen Höhen kann dies
vorkommen. Da es sich bei der Aorta um das größte Gefäß des
Körpers handelt ist die Prognose äußerst schlecht.

\paragraph{Nierenruptur}\index{Nierenruptur}\label{topic:nierenruptur}
Die Niere als naturgemäß gut durchblutetes Organ kann im
Falle einer Verletzung auch massive Probleme machen.


\subsection{Beckentrauma}

Ein Beckentrauma ist eine Verletzung des Beckens. In manchen Lehrbüchern werden sie zu den Bauchverletzungen gezählt. Besonders gefährlich sind die Beckenfrakturen, bei welchen binnen weniger Minuten bis zu fünf Liter Blut (also fast die gesamte Blutmenge) verloren gehen können. Es besteht höchste Schockgefahr!
\begin{itemize}
    \item Entstehung durch erhebliche Gewalteinwirkung  (Versch\"uttung, \"Uberrolltrauma)
    \item massive  Blutungen (bis 5\,l in wenigen  Minuten),
    Gefahr des Ausblutens (ev. ohne sichtbare  Blutungsquelle!)
    \item Mitverletzung von Organen im kleinen  Becken:
    \begin{itemize}
        \item Dickdarm ${\rightarrow}$ Blutung aus Anus
        \item Harnsystem ${\rightarrow}$ Blut im Harn, Harn in der
        Bauchh\"ohle
        %\item (Ausnahme: Spontanruptur der gef\"ullten 
        %Harnblase bei stumpfem Bauchtrauma)
        \item Hoden ${\rightarrow}$ H\"amatome am Skrotum und am  Damm
    \end{itemize}
\end{itemize}


\gMassnahmen{
    \paragraph{Spezielle Maßnahmen bei instabilen
    Beckentraumen}
    \begin{itemize}
        \item \gTxMassVitBes
        \begin{itemize}
            \item Lagerung: flach
            \item \ce{O2}: 10-15\,l
        \end{itemize}
        \item Schockbek\"ampfung
        \item Rettung mit Schaufeltrage oder KED-System
        \item Lagerung auf/Schienung mit Vakuummatratze (wenn Schocklagerung notwendig, Tragentisch schräg stellen)
        \item Bei Beckenfraktur:
        \index{Beckengurt}\gEs{Beckengurt}, zur Not mit
        G\"urtel
    \end{itemize}
}


%%%%%%%%%%%%%%%%%%%%%%%%%%%%%%%%%%%%%%%%%%%%%%%%%%%%%%%%%%%
%%%%%%%%%%%%%%%%%%%%%%%%%%%%%%%%%%%%%%%%%%%%%%%%%%%%%%%%%%%
%%%%%%%%%%%%%%%%%%%%%%%%%%%%%%%%%%%%%%%%%%%%%%%%%%%%%%%%%%%

\section{Polytrauma}

%\gAuthNotes{Neufassung}

\marginlabel{Polytrauma}\gDefinition{Gleichzeitige \gEs{Verletzung mehrerer K\"orperregionen} oder
Organsystemen, von denen \gEs{wenigstens eine} allein
\gEs{oder die Kombination mehrerer} \gEs{lebensgef\"ahrlich}
ist.}

Beispiele:
\begin{itemize}
    \item Bauchtrauma + Oberschenkel; SHT + WS + Thoraxtrauma etc.
    \item gebrochener Finger + gebrochene Zehe ${\rightarrow}$ kein
    Polytrauma
\end{itemize}

Häufigkeit verletzter Körperregionen bei Polytraumatisierten (nach \cite{BoehmerTaschRett6} und \cite{Gorgass7}):
\begin{itemize}
	\item 60-67 \% Schädel-Hirn
	\item 45-75 \% Extremitäten
	\item 25-30 \% Thorax
	\item 10-37 \% Abdomen und Becken
	\item 5-15 \% Wirbelsäule
\end{itemize}

\gFazit{Die Arbeitsdiagnose Polytrauma kann auch ausschließlich aufgrund des Unfallmechanismus und der Vitalparameter gestellt werden. Patient gilt als polytraumatisiert bis zum Beweis des Gegenteils! Jedes Polytrauma ist notarztpflichtig.}

Die Sterblichkeit bei polytraumatisierten Patienten liegt zwischen 20-50 \%, ist also hoch. Ein gezieltes, strukturiertes und zügiges Vorgehen ist daher unabdingbar! Im Idealfall sollten zwischen Eintreffen der Rettungskräfte und Transportbeginn nicht mehr als 30 Minuten vergehen. \cite{BoehmerTaschRett6}

\gMassnahmen{
    \paragraph{Spezielle Maßnahmen bei Polytrauma}
    \begin{itemize}
        \item Schneller Traumacheck gleichzeitig  mit  Sichern der Vitalparameter! (gründlicher Traumacheck erst später)
        \item Bewu{\ss}tsein
        \item Atmung
        	\begin{itemize}
			\item Atemwege freimachen und freihalten (ggf. Absaugen)
			\item \ce{O2}: 15\gLMin
			\item Atemwegssicherung durch Intubation (Notarzt!)
			\end{itemize}
        \item Kreislauf: 
			\begin{itemize}        	
        	\item Carotispuls tasten
        	\item periphere Pulse tasten (nur orientierend!)
        	\item Kapillardurchblutung überprüfen (Füllungszeit)
        	\item starke Blutungen stillen
        	\item Reanimationsbereitschaft und Monitoring (Verlaufskontrolle)
        	\item Wärmeerhaltung
        	\item venösen Zugang für Notarzt vorbereiten
        	\end{itemize}
        \item Lagerung entsprechend den Verletzungen (Vakuummatratze) und der Herz-Kreislauf-Situation -- Prioritäten absteigend:
        	\begin{itemize}
        	\item Atem- und Kreislaufstillstand $\rightarrow$ Rückenlage
			\item Bewusstlosigkeit, nicht intubiert $\rightarrow$ stabile Seitenlage (bei Thorax- und/oder Extremitätentrauma auf verletzte Seite legen)
			\item Atemnot bei Thoraxtrauma $\rightarrow$ Oberkörper hoch (45°)
			\item Abdominaltrauma $\rightarrow$ Schonhaltung, Knierolle
			\item SHT $\rightarrow$ Oberkörper hoch (30°)
			\item Wirbelsäulentrauma $\rightarrow$ Flachlagerung auf Vakuummatratze
			\item Extremitätentrauma $\rightarrow$ Extremität schienen und fixieren
			\item Schock $\rightarrow$ Schocklagerung durch schräg stellen des Tragentisches
			\end{itemize} 
        \item wenn einigerma{\ss}en stabil, gr\"undlicher  Traumacheck
        \item Verletzungsmuster/Unfallmechanismus erheben $\rightarrow$ zu erwartende Komplikationen einschätzen
        \item Notarzt: Schmerztherapie, Narkose, Intubation und Beatmung
        \item Transport:
        	\begin{itemize}
			\item Aviso im Zielspital (über Leitstelle)
			\item bei Wirbelsäulen-Trauma Rücksprache mit der Leitstelle wegen ev. Hubschrauber-Transport
			\end{itemize}
    \end{itemize}
}

% UNFALLORT IN DEN ABSCHNITT EINSATZTAKTIK VERSCHOBEN (Emhofer)

%\subsubsection{Der Unfallort und der Unfallmechanismus}

%\begin{itemize}
%    \item Absichern  der Unfallstelle,  m\"oglichst auff\"allig machen (Blaulicht,
% Warnblinkanlage, Scheinwerfer)
% \item bei mehreren Verletzten: Triage
% \item Nachalarmierung : andere  Rettungsmittel, Feuerwehr, Polizei
% \item Eingeklemmte im Fahrzeug stabilisieren,  nicht eigenm\"achtig bergen
% \item Crash -Rettung nur  bei  Fremdgef\"ahrdung, HWS stabilisieren!
%\end{itemize}

\paragraph{Polytrauma beim Kind} Der Schockindex ist beim Kind unbrauchbar, da die
Normalwerte andere sind! Kinder bleiben lange stabil, verfallen dann aber schlagartig!


%%%%%%%%%%%%%%%%%%%%%%%%%%%%%%%%%%%%%%%%%%%%%%%%%%%%%%%%%%%
%%%%%%%%%%%%%%%%%%%%%%%%%%%%%%%%%%%%%%%%%%%%%%%%%%%%%%%%%%%
%%%%%%%%%%%%%%%%%%%%%%%%%%%%%%%%%%%%%%%%%%%%%%%%%%%%%%%%%%%

\section{Extremitätentrauma und Sportverletzungen}
%Gelenke betreffend
Verletzungen der Extremitäten können in verschiedenen Formen auftreten:
\begin{itemize}
	\item Verstauchung (Distorsion)
	\item Verrenkung (Luxation)
	\item Knochenbruch (Fractur)
	\item Amputationen
\end{itemize}

\subsection{Verstauchung (Distorsion), Bänderzerrung und Bänderriss}
\gDefinition{Zerrung/\"Uberdehnung des Kapsel-/Bandapparates}

Bei der Verstauchung (\gT{Distorsion}) handelt es sich um eine Gelenksverletzung, bei welcher ein Knochen durch eine stumpfe Gewalteinwirkung kurz verschoben oder verdreht wird. Im Gegensatz zur Verrenkung springt der Knochen aber von selbst wieder in seine ursprüngliche Stellung im Gelenk zurück.

\paragraph{Symptome}
\begin{itemize}
    \item Schmerzen (vor allem Bewegungs- und Belastungsschmerz)
    \item Schwellung
    \item Blutergu{\ss} (H\"amatom)
    \item ev. Bewegungseinschränkungen
    \item hörbares Schnalzen beim Bänderriss
\end{itemize}

\gMassnahmen{
    \paragraph{Spezielle Maßnahmen bei einer Verstauchung}
    \begin{itemize}
		\item Kleidung und Schmuck von der betroffenen Gliedmaße entfernen ($\rightarrow$ Schwellung)
        \item ruhigstellen
        \item hochlagern
        \item k\"uhlen
        \item ad Unfallabteilung (mu{\ss} nicht sofort sein)
    \end{itemize}    
}

\subsection{Verrenkung (Luxation)}

\gDefinition{Teilweise oder vollst\"andige Verrenkung eines Gelenks}

Die Verrenkung (\gT{Luxation}) ist eine Gelenksverletzung, bei der -- im Gegensatz zur Verstauchung -- der verschobene oder verdrehte Knochen nicht von selbst in das Gelenk zurück gesprungen ist. Die betroffene Extremität kann daher nicht mehr bewegt werden. Bei einer Luxation dürfen keine eigenständigen Einrenkversuche unternommen werden!

\paragraph{Symptome}
\begin{itemize}
    \item Schmerzen
    \item Schwellung
    \item Blutergu{\ss}
    \item abnorme Gelenkstellung (federnd fixiert), bewegungsunf\"ahig
\end{itemize}

\gMassnahmen{
    \paragraph{Spezielle Maßnahmen bei einer Verrenkung}
	\begin{itemize}
		\item Kleidung und Schmuck von der betroffenen Gliedmaße entfernen ($\rightarrow$ Schwellung)
		\item Stellung der Gliedmaße belassen
		\item federnd fixierte Lagerung
		\item \gEs{keine Einrenkversuche!}
	\end{itemize}
}


%%%%%%%%%%%%%%%%%%%%%%%%%%%%%%%%%%%%%%%%%%%%%%%%%%%%%%%%%%%
%%%%%%%%%%%%%%%%%%%%%%%%%%%%%%%%%%%%%%%%%%%%%%%%%%%%%%%%%%%
%%%%%%%%%%%%%%%%%%%%%%%%%%%%%%%%%%%%%%%%%%%%%%%%%%%%%%%%%%%

\subsection{Knochenbr\"uche (Frakturen)}

Einen Knochenbruch nennt man auch \gT{Fraktur}. Man erkennt ihn am häufigsten durch eine Fehlstellung des Knochens und durch abnorme Beweglichkeit der betroffenen Gliedmaße. Seltener sind freie Knochenfragmente zu sehen. Je nach betroffenen Knochen kann es zu unterschiedlich großen Blutverlusten kommen:
\begin{itemize}
	\item \gEs{Oberarm: 100-800\,ml
	\item Unterarm: 50-400\,ml
%	\item Rippen: 2000-3000\,ml
	\item Becken: 500-5000\,ml (!)
	\item Oberschenkel: 300-2000\,ml
	\item Unterschenkel: 100-1000\,ml}
\end{itemize}
Die Blutung kann aus dem Knochenmark, zerrissener Muskulatur oder aus einem verletzten Gefäß kommen.

\subsubsection{Arten, Symptome und allgemeine Maßnahmen bei Knochenbrüchen}

\paragraph{Arten}
\begin{itemize}
    \item geschlossene Fraktur
    \item offene Fraktur (aufgrund der Gewalteinwirkung kommt es auch zu einer Wunde oberhalb der Bruchstelle)
\end{itemize}

% % gestrichen lt. STEUER 2009-08-29
% Weiter unterscheidet man:
% \begin{itemize}
%     \item komplette Fraktur (vollständige Durchtrennung)
%     \item inkomplette Fraktur (unvollständige Durchtrennung)
% \end{itemize}


\paragraph{Gefahren}
\begin{itemize}
    \item starke Blutungen, ev. bis Schock (z.B. Oberschenkelfraktur)
    \item begleitende Nervenverletzungen
    \item Fettembolie (Fraktur langer R\"ohrenknochen)
    \item Infektionsgefahr (Osteomyelitis bei offenen Frakturen)
    \item Kombination von mehreren Frakturen $\rightarrow$ großer Blutverlust
    \item schwere Muskelquetschung kann zu einem Sauerstoffmangel im Muskel führen (Kompartmentsyndrom)
\end{itemize}

Auch bei einer harmlos anmutenden Fraktur muss ein gründlicher Traumacheck durchgeführt werden, damit keine Begleitverletzungen übersehen werden. Dafür ist es unbedingt notwendig, dass der Patient entkleidet wird.\footnote{Nachdem der Patient für die Untersuchung entkleidet wurde, muss er nach der Untersuchung zugedeckt werden (Wärmeerhaltung)!} Beim Entfernen der Schuhe ist darauf zu achten, dass ein Sanitäter das Bein fixiert, während der andere Sanitäter mittels Stiefelgriff (siehe ...) vorsichtig den Schuh öffnet und schließlich von Fuß abzieht.

\paragraph{Frakturzeichen} \index{Frakturzeichen}

Es gibt sichere und unsichere Frakturzeichen:

\noindent\begin{tabularx}{\linewidth}[H]{XX}
\toprule
\multicolumn{2}{c}{\gEs{Frakturzeichen}}
\\\midrule
\multicolumn{1}{c}{\gEs{Sichere}}
 &
\multicolumn{1}{c}{\gEs{Unsichere}}
\\\toprule
\begin{itemize}
    \item Abnorme Beweglichkeit
    \emph{("`Pseudogelenk"')}\footnote{"`Pseudogelenk"':
    \emph{"`Wie ein Gelenk"'.} Beim Pesudogelenk hat es den
    Anschein, als würde ein Gelenk bestehen (welches nicht
    dorthin gehört).)}
    \item Fehlstellung
    \item Knochenreiben (Krepitus) \footnote{\so{NICHT}
    absichtlich ausprobieren!}
    \item Sichtbare freie Knochenteile (= offener Bruch)
\end{itemize}
 &
 \begin{itemize}
     \item Schmerz
     \item Schwellung
     \item H\"amatom
     \item gest\"orte Funktion / Bewegungsunf\"ahigkeit
 \end{itemize}
\\\bottomrule
\end{tabularx}
\captionof{table}{Frakturzeichen}

\begin{figure}[h]
    \begin{center}
\includegraphics[width=6cm]{./images/fraktur_fehlstellung_nativ.jpg}
\includegraphics[width=6cm]{./images/fraktur_fehlstellung_xray.jpg}
\caption{Fehlstellung -- In natura und in
Röntgendarstellung. \gSourcePicture{Hauer}\gAuthNotesPicLegalOk{}{}}
\end{center}
\end{figure}

\paragraph{Spezielle Maßnahmen}

Hauptaugenmerk bei der Versorgung der Knochenbrüche liegt auf der Blutstillung, der Schienung, sowie auf dem Erkennen von großen Blutverlusten.

\gMassnahmen{
    \paragraph{Spezielle Maßnahmen bei einer Fraktur}
	\begin{itemize}
		\item Kleidung und Schmuck von der betroffenen Gliedmaße entfernen ($\rightarrow$ Schwellung)
		\item Blutstillung, Wundversorgung
		\item bei starken Schmerzen oder großen Blutverlust $\rightarrow$ Notarzt
		\item Ruhigstellung/Schienung
		\item ggf. Schockbekämpfung
	\end{itemize}
}

\subsubsection{Spezielle Knochenbrüche}

\paragraph{Schlüsselbeinfraktur} Die Schlüsselbeinfraktur entsteht meistens, wenn der Patient mit ausgestreckten Armen stürzt. Die große Gefahr dieser Fraktur liegt darin, dass die Schlüsselbeinarterie und -vene verletzt werden können. Die Ruhigstellung erfolgt einerseits durch ein Armtragetuch (Dreiecktuch) und andererseits durch Fixation des Arms am Oberkörper durch ein zweites Dreiecktuch.

\paragraph{Schenkelhalsfraktur} Bei der Schenkelhalsfraktur handelt es sich um einen Bruch des Oberschenkelhalses. Meist sind ältere Menschen betroffen, die auf die Hüfte gestürzt sind. Symptome sind:
\begin{itemize}
	\item starker Bewegungs- oder Stauchungsschmerz
	\item betroffenes Bein ist nach außen gedreht und verkürzt
	\item Unfallhergang meist Sturz (auch an Kollaps denken!)
\end{itemize}
Die Versorgung erfolgt durch Rettung mittels Schaufeltrage und Ruhigstellung mittels Vakuummatratze. An große Blutverluste denken (Kontrolle der Vitalparameter)! Eventuell kann eine Schmerztherapie (Notarzt) notwendig sein.

\subsection{Amputationen}

\gDefinition{Amputation bezeichnet einen kompletten oder inkompletten Abriss bzw. Abtrennung (von Teilen) einer Extremität oder eines anderen Körperteils, sodass deren Durchblutung ganz oder teilweise aufgehoben ist. Den abgetrennten Körperteil nennt man Amputat. \cite{BoehmerTaschRett6}}

Amputationen erfolgen meist bei Motorradunfällen oder Schnittverletzungen (z.B. mit einer Kreissäge). Je glatter die Wundränder sind, desto größer ist die Chance, dass das Amputat wieder erfolgreich angenäht werden kann. Der Amputatsstumpf wird steril mit einem Druckverband versorgt (kein abbinden!). An große Blutungen und Schockgefahr denken!

\gMassnahmen{
    \paragraph{Versorgung eines Amputats}
	\begin{itemize}
		\item keine Reinigung des Amputats
		\item sterile Versorgung des Amputats (Verband)
		\item Lagerung und Transport in doppelwandigem Amputatsbeutel bei einer Kälte von ca. 4°C. Zwischen die beiden Schichten des Beutels gibt man Wasser mit Eis (im Verhältnis 1:1).
		\item Das Amputat darf auf keinen Fall gefroren werden oder direkten Kontakt zur kühlenden Substanz haben!
	\end{itemize}
}

\subsection{Weitere Verletzungsbilder}
Neben den bereits genannten Verletzungen, treten (im Sport oder in der Freizeit) oft auch folgende, meist harmlose, Verletzungen auf, welche hier der Vollständigkeit halber erwähnt sind:
\begin{itemize}
	\item Muskelkrampf
	\item Muskelzerrung
	\item Muskelfaserriss
	\item Achillessehnenriss
	\item Meniskusschaden
\end{itemize}

%%%%%%%%%%%%%%%%%%%%%%%%%%%%%%%%%%%%%%%%%%%%%%%%%%%%%%%%%%%
%%%%%%%%%%%%%%%%%%%%%%%%%%%%%%%%%%%%%%%%%%%%%%%%%%%%%%%%%%%
%%%%%%%%%%%%%%%%%%%%%%%%%%%%%%%%%%%%%%%%%%%%%%%%%%%%%%%%%%%
%%%%%%%%%%%%%%%%%%%%%%%%%%%%%%%%%%%%%%%%%%%%%%%%%%%%%%%%%%%

\section{Verbrennungen}

\gDefinition{Die Verbrennung ist eine durch Hitzeweinwirkung entstandene
Schädigung bzw. Verwundung der Haut.}

\marginlabel{\cite{Kamolz2009,Verbrennungsmed2009Praehosp,Beneker2004,Loennecker2001}}
Verbrennungen reichen von kleinen Brandblasen bis schweren Schädigungen der Haut und der darunter liegenden Gewebe (z.B. großflächige Verkohlung). Um die Schwere der Verletzung rasch einschätzen zu können muss die Ausdehnung und der Schweregrad der Verbrennung beurteilt werden. Dazu stehen die Prozentregel sowie die Definitionen der Schweregrade zur Verfügung.

Die Gr\"o{\ss}e des gesch\"adigten Areals wird in Prozent der
ver\-brannten K\"orperoberfl\"ache angegeben. Diese kann entweder mit der \gEs{Neunerregel} (siehe Tab. \ref{tab:Neunerregel} oder mit der \gEs{Handregel} ermittelt werden. Die Handregel besagt, dass die Handfläche des Patienten einem Prozent seiner Körperoberfläche entspricht.

\gMarried{
\begin{center}
\begin{tabulary}{\linewidth}[H]{LCC}
\toprule
 & \% &
\\\midrule
Kopf:
 & 9  &
\\
Arm, je:  &  9  &
\\
Rumpf vorne: & 18  &
\\
Rumpf hinten: & 18 &
\\
Bein, je: & 18 &
\\
Genitalien:  & 1 &
\\\bottomrule
\end{tabulary}
\end{center}
\caption{Neunerregel} \label{tab:Neunerregel}
}{
\includegraphics[width=\linewidth]{./images/noauth/AASdev20090228-img94.png}
}

Der Schweregrad einer Verbrennung wird wie folgt eingeteilt:
\begin{center}
\begin{tabulary}{\linewidth}[H]{CL}
\toprule
 Grad
 & Symptome
\\\midrule
1\,\textdegree
 &
R\"otung, Schmerz, Schwellung
\\
2\,\textdegree
 &
Blasenbildung, Schmerz
\\
3\,\textdegree
 &
\gT{Verschorfung}, Schmerzen
abhängig von der tiefe der Schädigung
(Zerstörung von Nerven), Schrumpf\-ung des Gewebes durch
Hitze\-ein\-wirkung, reicht bis zur Verkohlung \\\bottomrule
\end{tabulary}
\end{center}

Mit den erhobenen Werten (\% der verbrannten Körperoberfläche und Schweregrad) beurteilen Sie die Schockgefahr wie folgt:
\gCave{Ab \gEs{5\%} verbrannter (mind. 2. Grad) Körperoberfläche bei
Kindern oder \gEs{10\%} bei Erwachsenen besteht Schockgefahr
(Fl\"ussigkeitsverlust)!}


\gMassnahmen{
    \paragraph{Spezielle Ma{\ss}nahmen bei leichter
    Verbrennung}
    
    \begin{itemize}
        \item Eigenschutz!
        \item Kleiderbrand l\"oschen
        \item Kleidung rasch entfernen (sofern nicht  eingeschmolzen!)
        \item mind. 10 min. mit m\"a{\ss}ig kaltem Wasser die betroffenen Körperregionen spülen  --
       \gEs{NICHT unterk\"uhlen!}
        \item keimfreier, nicht verklebender Verband, locker anlegen
% TODO Stellungsnahme zu Water-Jel
%         \item wenn vorhanden: Water-Jel

        \item Schockbek\"ampfung
        \item NA wenn n\"otig
    \end{itemize}
    \gCave{CAVE Unterkühlung bei zu langer oder zu kalter
    Spülung!}
}

\gMassnahmen{
    \paragraph{Spezielle Ma{\ss}nahmen bei großflächiger
    Verbrennung}
    \begin{itemize}
        \item \gTxMassVitBes
        \begin{itemize}
            \item \ce{O2}: 10\gLMin
        \end{itemize}
        \item Kleidung rasch entfernen (sofern nicht  eingeschmolzen!)
        \item mind. 10 min. mit m\"a{\ss}ig kaltem Wasser die betroffenen Körperregionen spülen  --
       \gEs{NICHT unterk\"uhlen!}
        \item keimfreier, nicht verklebender Verband, locker anlegen
% TODO Stellungsnahme zu Water-Jel
%         \item wenn vorhanden: Water-Jel
%         \footnote{Streitfrage: Water-Jel: Umstritten}
        \item Schockbek\"ampfung
    \end{itemize}
    \gCave{CAVE Unterkühlung bei zu langer oder zu kalter
    Spülung!}
}



\paragraph{Was man besser \gE{nicht} tun sollte \ldots}

\begin{itemize}
    \item NEIN: Salbe, Puder, etc. auf frische  Brandwunde
    \item NEIN: Blasen \"offnen
    \item NEIN: festklebende Kleidung gewaltsam  entfernen
    \item NEIN: aluminiumbeschichteten Verband  verwenden (Gefahr der
    Hitzereflexion bei  nicht ausreichend gek\"uhlten  Verbrennungen!)
\end{itemize}



\subsection{Inhalationstrauma}

\gDefinition{Von einem Inhalationstrauma spricht man bei
Einatmung von reizenden, giftigen und/oder heißen Gasen,
welche eine Schädigung der Atemwege und der Lunge zur Folge
hat.}

\paragraph{Leitsymptome}

\begin{itemize}
    \item Z.n. Gasexposition
    \item angebrannte Gesichtshaare
    \item geschwollene, rußige Lippen
    \item Rußpartikel im Auswurf
    \item Husten
    \item Heiserkeit
    \item Stridor
    \item Atemnot
\end{itemize}

Das volle Ausmaß der Schädigung kann unmittelbar oder
innerhalb von Stunden erreicht werden. Eine engmaschige
Beobachtung des Patienten ist daher sehr wichtig.

\gMassnahmen{
    \paragraph{Spezielle Maßnahmen beim Inhalationstrauma}~

\begin{itemize}
    \item Vitale Bedrohung einschätzen. \par
    \gTxBeiVit \gTxMassVitBes
    \begin{itemize}
    \item \ce{O2}: 15\gLMin
\end{itemize}    
    \item Lagerung: Oberkörper hoch
     \item \ce{O2}: 6-8\gLMin
     \item Engmaschiges Monitoring
\end{itemize}

\ASBFLO{\gZ{Notfallsanitäter geben bei einer Rauch- oder
Reizgasinhalation ein inhalatives Kortikoid (z.B.
Pulmicort\texttrademark) laut Algorithmus der
Medikamentenliste 1.}}

}









%%%%%%%%%%%%%%%%%%%%%%%%%%%%%%%%%%%%%%%%%%%%%%%%%%%%%%%%%%%
%%%%%%%%%%%%%%%%%%%%%%%%%%%%%%%%%%%%%%%%%%%%%%%%%%%%%%%%%%%

\subsection{Verletzungen durch chemische Stoffe -- "`Chemische Verbrennungen"'}
\label{topic:veraetzung-trauma}\marginlabel{\cite{RD200811Hoffmann}}

Vergiftungen durch Einnahme von Säuren und Laugen werden im
Kapitel \emph{"`Krankheiten und Medizinische Notfälle"'}
unter \gSee{topic:veraetzung-intoxikation} besprochen. 

Schädigungen der haut durch Säuren, Laugen oder andere schädliche Stoffe weisen
-- trotz der unterschiedlichen Entstehung -- Gemeinsamkeiten mit
Verbrennungsschäden auf. Daher werden diese Verletzungen oft auch als
\emph{"`chemische Verbrennung"'} bezeichnet, auch wenn meist keine Hitze im
Spiel ist.  

Genau wie bei der Verbrennung kommt es zu einem Ausfall der Schutzfunktion der
Haut, oft zu Schmerzen und zu gefährlichen Schockzuständen. Auch die
Einschätzung der Ausdehnung der Verwundung erfolgt wie bei der Verbrennung mit
der Neuner-Regel o.ä. 

Eine Besonderheit stellt jedoch der verantwortliche Stoff dar, mit welchem u.U.
der Patient, dessen Kleidung und evt. auch die Umgebung kontaminiert ist. Von
diesem Stoff kann sowohl für den Patienten, als auch für den Helfer eine Gefahr
ausgehen. Auf den Eigenschutz und die Dekontamination muss großer Wert gelegt
werden. 

\gMassnahmen{
\paragraph{Spezielle Maßnahmen bei Verletzungen mit chemischen Substanzen}

\begin{enumerate}
    \item Auf Selbsschutz achten!
    \begin{itemize}
        \item Patienten ggfs. retten, aber nur wenn es der
        Eigenschutz zulässt.
    \end{itemize}
    \item Gefahr erkennen: \gEs{Welcher Stoff?}  Expertenrat einholen!
    \item Wenn erforderlich: Absperrmaßnahmen durchführen
    \item Wenn erforderlich:  Spezialkräfte anfordern.
    \item Vorgehen analog zu Abschnitt "`Gefahrengutunfall"'
    (\gSee{topic:gefahrengutunfall})
    \item Weitere Opfer ausschließen (lassen)
    \item Sicherung der Vitalfunktionen, Beurteilung ob der
    Patient vital bedroht ist.
    \item \gTxBeiVit \gTxMassVitBes
    \begin{itemize}
        \item Lagerung: situationsgerecht.
    \end{itemize}
    \item Entfernung kontaminierter Kleidung
    \item Sp\"ulung der betroffenen Hautpartien mit Wasser  (richtig abrinnen
lassen)
\item Auge: Sp\"ulung mit Wasser / Ringerlösung (nach au{\ss}en abrinnen
lassen!),  Kontaktlinsen entfernen
\item Hospitalisation: Ab einer Ausdehnung von 15\% (Erwachsener) bzw. 10\%
(Kind) Spezialabteilung für Verbrennungsverletzungen. Bei Augenverletzungen
Abt. für Augenheilkunde.
\end{enumerate}
}










%%%%%%%%%%%%%%%%%%%%%%%%%%%%%%%%%%%%%%%%%%%%%%%%%%%%%%%%%%%
%%%%%%%%%%%%%%%%%%%%%%%%%%%%%%%%%%%%%%%%%%%%%%%%%%%%%%%%%%%
%%%%%%%%%%%%%%%%%%%%%%%%%%%%%%%%%%%%%%%%%%%%%%%%%%%%%%%%%%%

\section{Erfrierungen}
\gDefinition{...sind Gewebsnekrosen, die durch lokale 
Durchblutungsst\"orungen aufgrund von  K\"alteeinwirkung (Eis, Wind, N\"asse,  K\"alte)
entstanden sind.}

Besonders gef\"ahrdet sind k\"orperferne K\"orperteile wie Zehen,
Finger, Ohren oder Nase.

Einteilung erfolgt in drei \gEs{Stadien} (siehe Tab. \ref{tab:Erfrierungen}), diese haben am Notfallort keine
entscheidende Bedeutung f\"ur die zu setzenden Ma{\ss}nahmen. Den
Schweregrad einer Erfrierung kann man erst nach Tagen voll
einsch\"atzen. Deshalb ist es wichtig, die Warnsymptome ernst zu
nehmen!

\begin{center}
\begin{tabulary}{\linewidth}[H]{CLC}
\toprule
 Stadien
 & Symptome
 &
\\\midrule
1
 &
Blässe, Empfindungsstörungen
 &
\\
2
 &
Blasenbildung, Schmerzen
 &
\\
3
 &
grau/schwarze Marmorierung (Nekrose), Gefühllosigkeit
\\\bottomrule
\end{tabulary}
\caption{Stadien der Erfrierung} \label{tab:Erfrierungen}
\end{center}


\gMassnahmen{
    \paragraph{Spezielle Ma{\ss}nahmen bei Erfrierungen}~
\begin{itemize}
\item beengende Kleidungsst\"ucke lockern (Einschn\"urung erh\"oht
Risiko einer Erfrierung)
\item Keimfreier, gepolsterter Verband
\item Warme, gezuckerte Getr\"anke -- KEIN Alkohol
\item K\"orperstamm w\"armen, NIE das erfrorene K\"orperteil abreiben!!
\end{itemize}
}




%%%%%%%%%%%%%%%%%%%%%%%%%%%%%%%%%%%%%%%%%%%%%%%%%%%%%%%%%%%
%%%%%%%%%%%%%%%%%%%%%%%%%%%%%%%%%%%%%%%%%%%%%%%%%%%%%%%%%%%
%%%%%%%%%%%%%%%%%%%%%%%%%%%%%%%%%%%%%%%%%%%%%%%%%%%%%%%%%%%

\section{Unfälle durch Strom- und Blitzschlag}

\marginpar{\includegraphics[width=\linewidth]{./images/electricity_japan.jpg}
\small\emph{Foto: Fina}}

Stromunfälle kommen typischerweise in folgenden Situationen vor:
\begin{itemize}
	\item Spielende Jugendliche oder Arbeiter auf Eisenbahnwagons
	\item LKW-Hebekräne oder Bagger geraten in Freileitungen
	\item Unachtsamer Umgang mit Elektrogeräten im Badezimmer
	\item Unsachgemäße Reparatur von Elektrogeräten
	\item Arbeitsunfälle (Elektriker, Heimwerker)
	\item Tragen von  Metallteilen in der Nähe von Hochspannungsleitungen
\end{itemize}

\subsubsection{Allgemeines zum elektrischen Strom und seine
Wirkung auf den Körper} 

Die Gefährdung bei Stromunfällen
hängt von verschiedenen Faktoren ab:
\begin{itemize}
	\item \gEs{Elektrische Spannung:} Die elektrische Spannung
	ist die Ursache für den elektrischen Stromfluss und wird
	in der Einheit Volt (V) gemessen. Wie gefährlich eine Spannung ist hängt nicht nur von der Höhe der Spannung selbst, sondern wesentlich auch vom elektrischen Strom, der von ihr hervorgerufen wird, ab. So kann z.B. die 12 V-Spannung einer Modelleisenbahn völlig ungefährlich sein, gibt es jedoch im Auto bei einer (großen) 12 V-Batterie einen Kurzschluss, kann ein gefährlich großer Strom fließen!
		\begin{itemize}
		\item \gT{Niederspannung}: Von Niederspannung spricht man bei Spannungen bis 1000 V. An einer Autobatterie liegt eine Spannung von 12 V. An einer Steckdose liegt eine Spannung von 230 V, in Industriebetrieben wird zusätzlich oft eine Spannung von 400 V verwendet. Diese Spannungen sind auf jeden Fall lebensgefährlich. \gZ{Ein Defibrillator schockt mit etwa 1000 V (!), ein externer Herzschrittmacher arbeitet mit bis zu 400 V.}
		\item \gT{Hochspannung}: Hochspannung ist eine Spannung über 1000 V. \gZ{Über Hochspannungsleitungen wird die elektrische Energie mit einer Spannung von 110 000 V (= 110 kV) bis 380 000 V (= 380 kV) übertragen. An den Oberleitungen der Eisenbahn liegen 15 000 V (= 15 kV). Bei Gewitterblitzen treten Spannungen von 10-30 Millionen Volt (10-30 Megavolt) auf.}
		\end{itemize}
	\item \gEs{Elektrische Stromstärke:} Elektrischer Strom ist der Fluss von Elektronen. Die treibende Kraft hinter dem Stromfluss ist die elektrische Spannung. Der elektrische Strom wird durch die elektrische Stromstärke beschrieben, welche mit der Einheit Ampere (A) gemessen wird. Kommt es zu einem Stromunfall ist die Schädigung im Körper von der Höhe der Stromstärke abhängig \footnote{Die Werte beziehen sich auf Wechselstrom, die Grenzwerte für Gleichstrom sind höher. Da oft die Stromart nicht bekannt ist, ist es prinzipiell sinnvoll von der gefährlicheren Stromart auszugehen!}:
\gZ{
		\begin{itemize}
		\item Bei geringen Stromstärken \gEs{unter 1 Milliampere (mA)} ist nur ein elektrischer Schlag spürbar.
		\item Bei Stromstärken \gEs{ab ca. 10 Milliampere (mA)} kommt es zu Schmerzempfindungen und Muskelkontraktionen. \cite{Jaros1}
		\item Bei Stromstärken \gEs{ab 15-25 mA} kommt es aufgrund der Muskelkontraktionen zu krampfhaften Festhalten der berührten Stromquelle. \cite{BoehmerTaschRett6}
		\item Stromstärken \gEs{ab 25-50 mA} sind \gEs{lebensgefährlich}! \cite{BoehmerTaschRett6}
		\item Stromstärken \gEs{ab 80 mA} führen zu Atemlähmungen und Kammerflimmern! \cite{Jaros1}
		\item Stromstärken \gEs{ab 3 A} führen zusätzlich zu Verbrennungen, Verkochungen und Verkohlungen. \cite{LPN3}
		\end{itemize}
}
	\item \gEs{Elektrischer Widerstand:} Jeder Körper und
	jeder Stoff setzt dem elektrischen Stromfluss einen
	Widerstand entgegen. Dieser elektrische Widerstand wird in
	Ohm ($\Omega$) gemessen. Je größer der Widerstand eines
	Körpers, desto kleiner ist der elektrische Strom. Trockene
	Haut hat einen elektrischen Widerstand von 100 000 Ohm,
	\gEs{feuchte Haut hat einen geringeren Widerstand} von
	weniger als 1000 Ohm. Der Widerstand von geschädigter
	(z.B. verbrannter) Haut ist noch geringer. \gZ{Greift man z.B. mit feuchten Händen in die Steckdose (230 V) so stellt sich durch den Körperwiderstand von ca. 1000 Ohm ein elektrischer Strom der Stärke 230 mA ein.}\footnote{Der Leser kann dies durch das Ohmsche Gesetz leicht selbst nachrechnen! Die Spannung $U$ ist das Produkt aus dem Widerstand $R$ und der elektrischen Stromstärke $I$: $U=R \cdot I$} Diese Stromstärke ist absolut lebensgefährlich, wie man aus der Übersicht oben schnell entnehmen kann.
	\item \gEs{Einwirkdauer:} Je länger die Einwirkzeit (Kontaktzeit), desto größer ist die Schädigung.
	\item \gEs{Stromweg:} Die Schädigung im Körper hängt auch vom Weg des Stromes durch den Körper ab. Fließt der Strom z.B. von einer Hand zur anderen Hand durch den Körper ist dies gefährlicher, als wenn der Strom von einer Hand zum Fuß fließt.
	\item \gEs{Stromart:} Wechselstrom ist gefährlicher als Gleichstrom. An den Steckdosen im Haushalt liegt Wechselspannung an. Dieser ändert die Stromrichtung 50 mal pro Sekunde (\gZ{Frequenz = 50 Hz}), weshalb die Wahrscheinlichkeit für Kammerflimmern auch bei kurzen Stromkontakten steigt. \cite{RedelsteinerHRNS1} \footnote{Bei monophasischen Defibrillatoren wird mit Gleichstrom gearbeitet, biphasische Defibrillatoren polen die Stromrichtung einmal um (Wechselstrom).} 
\end{itemize}

\gAuthNotes{\#9: EMHOFER: Bilder von verschiedenen
Steckdosen, Kabel und Hinweisschildern einfügen}

%\gT{Niederspannung}: bis 1000\,Volt (Steckdose:  230 Volt)
%\gT{Hochspannung}: \"uber 1000\,Volt, Stromst\"arke
%{\textgreater}  3 Ampere (\"Uberlandleitungen) 

\subsubsection{Besonderer Gefahrenbereich}
\index{Gefahrenbereich!Strom}\label{topic:gefahrenbereich-strom}
Starkstrom/Hochspannung durch Fachmann (E-Werk, ÖBB,
Betriebsleitung, ...) abschalten bzw. erden lassen 

Strom
gegen Wiedereinschalten sichern! Patient mit isolierendem
Material aus dem Stromkreis befreien. Sicherheitsabstand zu
stromführenden Freileitungen einhalten! 

\subparagraph{Freileitungen}


Bei Freileitungen mit
220 000 Volt (220 kV) sind mindestens 4 m Abstand zu halten.
\cite{Jaros1} 

Bei herabhängenden Hochspannungsleitungen oder
nach Blitzeinschlägen auf die Schrittspannung achten: Die
Spannung fällt im Boden nicht sofort auf 0 Volt ab, innerhalb
eines gewissen Umkreises kann daher zwischen den beiden
Beinen am Boden eine so genannte Schrittspannung auftreten.
Befinden Sie sich innerhalb eines solchen Gebietes lassen Sie
Ihre Beine eng beisammen und springen aus diesem Bereich
heraus. Je nach Bodenbeschaffenheit können Schrittspannungen
bis zu einer Distanz von 20 m von der herunterhängenden
Hochspannungsleitung auftreten. \cite{RedelsteinerHRNS1} 

\subparagraph{Bahnanlagen}
Bei
Notfällen auf Bahngleisen ist auf jeden Fall der
Streckenabschnitt von der ÖBB sperren zu lassen (über die
Leitstelle)! An den Masten der Oberleitungen der ÖBB sind
entsprechende Mastennummer angebracht. Diese müssen bei der Leitstelle gemeldet werden.
 
\subsubsection{Patientenversorgung}

\paragraph{Symptome} Die Verletzungsmuster und Symptome variieren je nach Spannung (Volt) und Stromst\"arke (Ampere):

\gParallel{
\begin{itemize}
	\item Unfallmechanismus und Unfallhergang
    \item Reizeffekte (Muskel-/ Nervenerregung):
    \begin{itemize}
        \item Verkrampfung ({\quotedblbase}kann den stromführenden Teil nicht
        loslassen{\textquotedblleft})
        \item Sensibilit\"atsst\"orungen
        \item ZNS ${\rightarrow}$ Bewu{\ss}tseinsst\"orung
        \item Herzrhythmusst\"orungen (Kammerflimmern, ev.  als Sp\"atfolge!)
    \end{itemize}
    \item Bewusstseinsstörungen bis Krampfanfälle oder Bewusstlosigkeit
    \item Lähmungen (vorübergehend)
    \item W\"armeentstehung
    \begin{itemize}
        \item Strommarken (Ein-/ Austrittsverbrennung)
        \item Lichtbogen: oberfl\"achliche und tiefe  Verbrennungen
        \item Blitzfiguren an der Haut
    \end{itemize}
\end{itemize}

%\begin{itemize}
%    \item Gleichstrom
%    \item Wechselstrom ${\rightarrow}$ eher  Herzrhythmusst\"orungen
%    \item Patient  kann Spannungsquelle ev. nicht  loslassen (Krampf der
%    Handmuskulatur)
%    \item Auch Defi stellt Strom-Gefahrenquelle  dar!!!
%    \item Blitzunfall
%    \begin{itemize}
%        \item mehrere Millionen (!) Volt
%        \item bis 100.000 Ampere
%    \end{itemize}
%\end{itemize}

%\paragraph{Symptome}
%\begin{itemize}
%    \item Bewu{\ss}tlosigkeit
%    \item L\"ahmungen (vor\"ubergehend)
%    \item Blitzfiguren auf der Haut
%    \item Herzrhythmusst\"orungen m\"oglich
%\end{itemize}
}{
\gCave{Das eigentliche Ausma{\ss} der Verletzung ist oft
{\quotedblbase}unsichtbar{\textquotedblleft}, da der Strom im
K\"orperinneren -- also {\quotedblbase}unterirdisch{\textquotedblleft}
-- viel Schaden anrichtet. }
}


\paragraph{Blitzschlag} Bei Gewitterblitzen treten
Spannungen von Millionen von Volt und Stromstärken von bis
zu 100\,000 Ampere auf. Trotzdem gibt es auch bei Blitzunfällen immer wieder Überlebende, da die Einwirkdauer extrem kurz (wenige Millisekunden) ist. Typische Symptome sind
\begin{itemize}
	\item Bewusstlosigkeit
	\item vorübergehende Lähmung,
	\item Blitzfiguren auf der Haut.
	\item Herzrhythmusstörugnen sind möglich
\end{itemize}


\gMassnahmen{
    \paragraph{Spezielle Ma{\ss}nahmen bei Stromunfällen}
    \begin{itemize}
        \item \gEs{Eigenschutz!} Bevor Sie den Patienten berühren, vergewissern Sie sich, dass der Patient keinen Kontakt mehr mit dem Stromkreis hat! \gEs{Schalten Sie auf jeden Fall vor jeder Berührung den Strom ab!}
        \begin{itemize} 
            \item Gefahrenbreich beachten:
            \gSee{topic:gefahrenbereich-strom}
                        \item Sicherheitsabstand
            (\gSee{topic:gefahrenbereich-strom})
            \item Schrittspannung und Spannungskegel
            beachten \gSee{topic:gefahrenbereich-strom}
            \item Bahnanlagen sperren/sichern lassen
            (\gSee{topic:gefahrenbereich-strom})
            \item Schalter ausschalten, (Haupt-)Sicherung
            auslösen, Starkstrom/Hochspannung durch Fachmann abschalten bzw. erden lassen
            \item Strom gegen Wiedereinschalten sichern!
            \item Patient mit isolierendem Material aus dem Stromkreis befreien

        \end{itemize}
        \item Beurteilung der vitalen Bedrohung.
        \par \gTxBeiVit
        \gTxMassVit
                \item \gEs{Monitoring \&
                Reanimationsbereitschaft}: auf evt.
                Herzrhythmusst\"orungen vorbereitet sein
                (bis zu 24\,h Latenz)
        \item Genaue Patientenuntersuchung ${\rightarrow}$
        Eintritts- und Austrittsmarke
        \item Hospitalisierung auch bei symptomlosen
        Patienten zur Überwachung
        \item Stromverbrennungen wie Verbrennungen versorgen
    \end{itemize}
}



\nocite{PHTLS6}
\nocite{emergency9}
\nocite{OHCSP7}
\nocite{Helfen2008}
\nocite{DRAn3}
\nocite{Patzelt2005}
\nocite{Ficklscherer2005}
\nocite{Sachsenweger2003}
\nocite{SiewertChirurgie8}




\gChapterBibliography
